% ============================================================
%  ROZDZIAŁ 3 — PROBLEM I SPOSÓB JEGO ROZWIĄZANIA
%  Plik: 03_problemiarchitektura.tex
%  Włączany do Thesis.tex przez % ============================================================
%  ROZDZIAŁ 3 — PROBLEM I SPOSÓB JEGO ROZWIĄZANIA
%  Plik: 03_problemiarchitektura.tex
%  Włączany do Thesis.tex przez % ============================================================
%  ROZDZIAŁ 3 — PROBLEM I SPOSÓB JEGO ROZWIĄZANIA
%  Plik: 03_problemiarchitektura.tex
%  Włączany do Thesis.tex przez % ============================================================
%  ROZDZIAŁ 3 — PROBLEM I SPOSÓB JEGO ROZWIĄZANIA
%  Plik: 03_problemiarchitektura.tex
%  Włączany do Thesis.tex przez \input{03_problemiarchitektura}
% ============================================================

\chapter{Problem i~sposób jego rozwiązania}

Niniejszy rozdział precyzuje formalną definicję problemu broken lineage
i~jego ujęcie jako zadania klasyfikacji binarnej (sekcja~\ref{sec:sformulowanie}).
Sekcja~\ref{sec:architektura} opisuje architekturę zaproponowanego
potoku przetwarzania, a~sekcja~\ref{sec:dane} charakteryzuje zastosowany
zbiór danych. Metodykę symulacji brakujących krawędzi przedstawiono
w~sekcji~\ref{sec:scenariusze}, natomiast szczegóły implementacji
algorytmów omówiono w~sekcji~\ref{sec:algorytmy}.

% ============================================================
\section{Sformułowanie problemu}
\label{sec:sformulowanie}
% ============================================================

Graf lineage modeluje się jako graf skierowany $G^{*} = (V, E^{*})$,
gdzie $V$ oznacza zbiór obiektów bazodanowych (tabel, widoków, zadań ETL),
a~$E^{*}$ — zbiór wszystkich rzeczywistych relacji przepływu danych
(\emph{data\_flow}). W~praktyce część krawędzi jest niedostępna
z~powodu zjawiska broken lineage: tabele tymczasowe nie~są rejestrowane
w~standardowych słownikach bazy danych, a~funkcje definiowane przez
użytkownika (\emph{User-Defined Function}, UDF) ukrywają wewnętrzną
logikę transformacji. Obserwowany graf ma~zatem postać
$G_{\mathrm{obs}} = (V, E_{\mathrm{obs}})$, gdzie
$E_{\mathrm{obs}} \subsetneq E^{*}$. Brakujące krawędzie
$E^{*} \setminus E_{\mathrm{obs}}$ reprezentują zerwane zależności,
których odtworzenie jest celem niniejszej pracy.

Zadanie wykrywania brakujących krawędzi sprowadza się do~klasyfikacji
binarnej. Dla~każdej pary węzłów $(u, v) \in V \times V$ algorytm
wyznacza wartość $\hat{y}(u, v) \in [0, 1]$, interpretowaną jako
oszacowane prawdopodobieństwo istnienia krawędzi w~pełnym grafie
$G^{*}$. Para jest klasyfikowana jako krawędź istniejąca (klasa~1)
lub~nieistniejąca (klasa~0) po~przekroczeniu progu binaryzacji
wynoszącego $0{,}5$.

Ze~względu na~strukturę grafu DLG-DG-23 kandydatami do~wykrywania
są wyłącznie pary węzłów różnych typów: (\emph{job},~\emph{table})
lub~(\emph{table},~\emph{job}). Para (\emph{table},~\emph{table})
nigdy nie~jest połączona krawędzią \emph{data\_flow}, a~relacja
\emph{parent\_child} łączy wyłącznie węzły \emph{table} i~\emph{field}
— pary te nie~są przedmiotem wykrywania. Ograniczenie przestrzeni
kandydatów do~typologicznie dopuszczalnych par eliminuje trywialnie
fałszywe przykłady ujemne i~zapewnia spójność z~semantyką domeny.

% ============================================================
\section{Architektura rozwiązania}
\label{sec:architektura}
% ============================================================

Potok przetwarzania składa się z~czterech głównych komponentów:
modułu wczytywania danych (\emph{loader}), modułu podziału krawędzi
(\emph{splitter}), właściwych algorytmów wykrywania oraz modułu
ewaluacji.

% TODO (WYMAGANE): Wstawić rysunek blokowy potoku:
%   JSON → loader → NetworkX DiGraph → splitter → algorytmy → metryki.
%   Rysunek musi być wprowadzony zdaniem przed jego wystąpieniem,
%   np. „Rysunek~X przedstawia schemat blokowy zaproponowanego potoku."

\subsection{Moduł wczytywania danych}

Dane wejściowe dla~każdego grafu DLG przechowywane są w~dwóch plikach
JSON: \texttt{DLGi-node.json} zawiera opisy węzłów wraz z~atrybutem
\texttt{asset\_type}, natomiast \texttt{DLGi-edge.json} — opisy
krawędzi z~atrybutami \texttt{relation\_type} i~\texttt{relation\_id}.
Moduł \texttt{loader.py} wczytuje oba pliki dla~każdego grafu i~buduje
obiekt \texttt{nx.DiGraph} z~biblioteki NetworkX, a~następnie zwraca
słownik wszystkich 18~grafów zbioru DLG-DG-23. Sam graf przechowuje
atrybut \texttt{name} (np.\ \texttt{DLG1}).

\subsection{Moduł podziału krawędzi}

Moduł \texttt{splitter.py} wyodrębnia ze~struktury grafu krawędzie
typu \emph{data\_flow} i~dzieli je~na~trzy rozłączne podzbiory
w~proporcji 60/20/20 (treningowy, walidacyjny, testowy). Próbki
pozytywne stanowią krawędzie usunięte ze~zbioru treningowego
i~przeznaczone do~ewaluacji; próbki negatywne to~pary węzłów
typologicznie dopuszczalnych (wyłącznie \emph{job}--\emph{table}),
które nigdy nie~były połączone krawędzią. Stosunek próbek negatywnych
do~pozytywnych wynosi 1:1, co~zapewnia zbalansowany zbiór ewaluacyjny
i~stabilne wartości metryk porównawczych.

Moduł \texttt{scenario\_splitter.py} realizuje analogiczny podział,
lecz zamiast losowego usuwania krawędzi stosuje trzy scenariusze
symulacji broken lineage opisane w~sekcji~\ref{sec:scenariusze}.

\subsection{Moduł ewaluacji}

Moduł \texttt{metrics.py} wyznacza standardowe metryki klasyfikacji
binarnej: precyzję, czułość, miarę F1, AUC-ROC i~AUC-PR. Próg
binaryzacji wynosi $0{,}5$ dla~wszystkich metod. Wyniki są
uśredniane po~wszystkich 18~grafach zbioru DLG-DG-23 i~agregowane
według scenariusza oraz algorytmu, co~umożliwia bezpośrednie
porównanie między metodami.

% ============================================================
\section{Zbiór danych DLG-DG-23}
\label{sec:dane}
% ============================================================

Zastosowany zbiór danych pochodzi z~pracy Chen et~al.~\cite{Chen2024}
i~zawiera 18~grafów lineage pobranych ze~środowiska produkcyjnego
Huawei Cloud. Dane zostały zanonimizowane — nazwy tabel i~atrybutów
zastąpiono identyfikatorami w~formacie UUID, a~schematy bazy danych
nie~zostały udostępnione. Zbiór opublikowano w~2024~roku jako pierwszy
otwarty zestaw rzeczywistych grafów lineage; szczegółową charakterystykę
poszczególnych grafów zamieszczono w~tabeli~\ref{tab:dlg-stats} w~sekcji~\ref{sec:dane-opis}.

Zbiór DLG-DG-23 wybrano ze~względu na~jego unikatowy charakter:
jest to~jedyny publicznie dostępny zestaw rzeczywistych grafów lineage,
co~umożliwia ewaluację algorytmów na~danych o~autentycznej topologii
przemysłowej, bez~ryzyka artefaktów wprowadzonych przez dane
syntetyczne. Anonimizacja nazw węzłów ma~bezpośredni wpływ na~dobór
algorytmów: metryki podobieństwa oparte na~nazwach obiektów
(np.~odległość Jaro--Winkler stosowana przez Boińskiego
et~al.~\cite{Boinski2025}) są nieaplikowalne. W~niniejszej pracy
wektor cech ograniczono zatem do~atrybutów topologicznych grafu,
co~jednocześnie stanowi element badań nad~pytaniem, czy sama struktura
grafu lineage jest wystarczająca do~wykrywania zerwanych zależności.

\subsection{Charakterystyka grafów}
\label{sec:dane-opis}

Tabela~\ref{tab:dlg-stats} przedstawia statystyki poszczególnych
grafów zbioru DLG-DG-23.

\begin{table}[ht]
\centering
\caption{Charakterystyka grafów zbioru DLG-DG-23.}
\label{tab:dlg-stats}
\begin{threeparttable}
\begin{tabular}{llrl}
\toprule
\textbf{ID} & \textbf{Węzły} & \textbf{Krawędzie DATA\_FLOW} & \textbf{Skala} \\
\midrule
DLG1  &    298 &    24 & mały   \\
DLG2  &    464 &    37 & mały   \\
DLG3  &    603 &    57 & mały   \\
DLG4  &    415 &    32 & mały   \\
DLG5  & 1\,299 &   175 & średni \\
DLG6  & 13\,840 & 1\,443 & duży  \\
DLG7  &    574 &    43 & mały   \\
DLG8  &    546 &    34 & mały   \\
DLG9  &    756 &    62 & mały   \\
DLG10 & 5\,378 &   497 & duży   \\
DLG11 & 2\,024 &   121 & średni \\
DLG12 & 5\,645 &   469 & duży   \\
DLG13 & 2\,157 &   229 & średni \\
DLG14 &    453 &    24 & mały   \\
DLG15 &    558 &    41 & mały   \\
DLG16 & 5\,574 &   745 & duży   \\
DLG17 &    651 &    38 & mały   \\
DLG18 & 17\,085 & 1\,694 & duży \\
\bottomrule
\end{tabular}
\begin{tablenotes}
\small
\item Skala: mały $<$ 1\,000 węzłów; średni 1\,000–5\,000; duży $>$ 5\,000.
\item Liczba krawędzi DATA\_FLOW podana dla sumy zbiorów treningowego,
      walidacyjnego i~testowego (podział 60/20/20).
\end{tablenotes}
\end{threeparttable}
\end{table}

% ============================================================
\section{Metodyka symulacji broken lineage}
\label{sec:scenariusze}
% ============================================================

Ponieważ zbiór DLG-DG-23 zawiera w~pełni zaobserwowane grafy lineage,
zjawisko broken lineage symulowane jest poprzez kontrolowane usuwanie
krawędzi. Podejście to jest standardowe w~literaturze wykrywania
krawędzi~\cite{LibenNowell2007}: dysponując kompletnym grafem
$G^{*}$, usuwa się podzbiór krawędzi, tworząc graf obserwowany
$G_{\mathrm{obs}}$, a~następnie sprawdza, czy algorytm potrafi
odtworzyć usunięte krawędzie. W~niniejszej pracy przyjęto trzy
scenariusze symulacji, odzwierciedlające typowe przyczyny powstawania
broken lineage.

\textbf{Scenariusz A} — losowe usunięcie krawędzi — realizuje
najprostszą formę symulacji: 20\% krawędzi \emph{data\_flow} jest
usuwanych ze~zbioru treningowego w~sposób losowy, bez~względu na~typ
węzłów końcowych. Scenariusz ten modeluje sytuację, w~której~broken
lineage wynika z~przyczyn losowych (np.~niedokompletowania logów
systemowych) i~nie~jest skoncentrowany na~konkretnym typie obiektu.

\textbf{Scenariusz B} — usunięcie krawędzi tabeli pośredniej —
symuluje sytuację, w~której tabela pośrednia (staging) jest tabelą
tymczasową i~nie~zostaje zarejestrowana w~metadanych. Usuwane są
wszystkie krawędzie prowadzące do~wybranego węzła \emph{table}
o~wysokim stopniu (tabela pośrednia), co~powoduje odcięcie
całego podgrafu zależności przepływającego przez ten węzeł.

\textbf{Scenariusz C} — usunięcie krawędzi do~tabeli docelowej —
symuluje ukrycie zależności na~ostatnim etapie potoku ETL, kiedy
krawędź do~tabeli docelowej (\emph{data mart}) nie~zostaje utrwalona.
Usuwane są krawędzie wchodzące do~wybranego węzła \emph{table}
o~zerowym lub~niskim stopniu wyjściowym, co~odpowiada tabelom
końcowym w~hierarchii przepływu danych.

% TODO (WYMAGANE): Wstawić tabelę zbiorczą scenariuszy:
%   Scenariusz | Co usuwamy | Co symuluje | Typ węzłów dotkniętych
%   Tabela musi być poprzedzona zdaniem „Tabela X przedstawia...".

% TODO (OPCJONALNIE): Wstawić poglądowe rysunki małego grafu
%   przed i~po usunięciu krawędzi dla~każdego scenariusza.
%   Rysunki muszą być wprowadzone zdaniem w~tekście.

% ============================================================
\section{Zaimplementowane algorytmy}
\label{sec:algorytmy}
% ============================================================

W~pracy zaimplementowano trzy grupy algorytmów wykrywania
brakujących krawędzi. Pierwszą grupę stanowią heurystyki grafowe
— metody niewymagające uczenia, opierające się wyłącznie na~strukturze
grafu. Drugą grupę tworzą klasyczne metody uczenia maszynowego,
budujące klasyfikator na~podstawie wektora cech topologicznych
pary węzłów. Trzecia grupa obejmuje metody zanurzeń sieciowych
(\emph{network embeddings}), reprezentowane przez algorytm Node2Vec.

% ============================================================
\subsection{Heurystyki grafowe}
% ============================================================

Heurystyki grafowe stanowią najprostszą klasę metod wykrywania krawędzi
(\emph{link prediction}). Wynik dla~pary węzłów $(u, v)$ wyznaczany jest
wyłącznie na~podstawie lokalnej lub~globalnej struktury grafu, bez~etapu
trenowania modelu. Metody te charakteryzują się niską złożonością
obliczeniową i~pełną interpretowalnością, co~czyni je~naturalnym punktem
odniesienia (ang.\ \emph{baseline}) przy~porównywaniu bardziej
zaawansowanych podejść~\cite{LibenNowell2007}.

Wstępna analiza struktury grafów DLG-DG-23 wykazała, że~miary oparte
na~wspólnych sąsiadach — Common Neighbors, indeks Jaccarda
i~miara Adamic-Adar — są nieskuteczne w~grafach dwudzielnych.
W~strukturze bipartite węzły tego samego typu (\emph{job}--\emph{job}
lub~\emph{table}--\emph{table}) nigdy nie~posiadają wspólnych sąsiadów
bezpośrednich, ponieważ krawędzie \emph{data\_flow} łączą wyłącznie
węzły różnych typów. W~konsekwencji wymienione miary zwracają wynik
zerowy dla~każdej typologicznie dopuszczalnej pary kandydackiej,
co~dyskwalifikuje je~jako metody wykrywania w~omawianym scenariuszu.
W~implementacji uwzględniono wyłącznie heurystyki działające poprawnie
na~grafach dwudzielnych: Preferential Attachment, L3, Katz
oraz~Personalized PageRank.

\subsubsection{Preferential Attachment}

Miara \emph{Preferential Attachment} (PA) opiera się na~obserwacji,
że~w~sieciach rzeczywistych nowe połączenia powstają z~prawdopodobieństwem
proporcjonalnym do~stopnia istniejących węzłów — węzły o~wielu sąsiadach
z~większym prawdopodobieństwem pozyskują kolejne
połączenia~\cite{Barabasi1999}. Jako heurystyka wykrywania krawędzi PA
została zastosowana przez Liben-Nowella i~Kleinberga~\cite{LibenNowell2007}.

Dla~pary węzłów $(u, v)$ wynik PA definiuje się jako:
%
\[
  \mathrm{PA}(u, v) = |N(u)| \cdot |N(v)|,
\]
%
gdzie $N(u)$ oraz $N(v)$ oznaczają zbiory sąsiadów odpowiednio węzła $u$
i~węzła $v$, a~$|\cdot|$ oznacza rozmiar zbioru, tj.\ stopień węzła w~grafie.

W~kontekście grafów lineage wysoki stopień węzła \emph{job} oznacza,
że~dane zadanie czyta z~wielu tabel źródłowych lub~zapisuje do~wielu
tabel docelowych; wysoki stopień węzła \emph{table} wskazuje na~tabelę
centralną — często pośrednią warstwę integracyjną w~potoku ETL.
Para $(u, v)$, w~której oba węzły mają duże stopnie, otrzymuje wysoki
wynik PA, co~interpretuje się jako sygnał potencjalnej zależności
\emph{data\_flow}. Miara PA nie~wymaga istnienia wspólnych sąsiadów,
co~stanowi jej kluczową przewagę w~grafach dwudzielnych.

\subsubsection{L3}

Miara L3 (ang.\ \emph{path of length 3}) liczy ścieżki długości~3
między węzłami $u$ i~$v$~\cite{LibenNowell2007}. W~grafie dwudzielnym
ścieżka długości~3 między węzłem \emph{job} a~węzłem \emph{table}
przebiega przez pośredni węzeł \emph{table} i~pośredni węzeł \emph{job}
— co odpowiada wzorcowi: zadanie~$u$ czyta z~tabeli~$t_1$, tabela~$t_1$
jest zapisywana przez zadanie~$j_1$, zadanie~$j_1$ czyta z~tabeli~$v$.
Wysoka wartość L3 sygnalizuje, że~para $(u, v)$ jest gęsto połączona
przez pośrednie obiekty, co~w~kontekście broken lineage może wskazywać
na~zależność ukrytą przez tabelę tymczasową.

Wynik L3 dla~pary $(u, v)$ definiuje się jako:
%
\[
  \mathrm{L3}(u, v) = \sum_{w \in V} A_{uw} \cdot A_{wv}^{(2)},
\]
%
gdzie $A$ oznacza macierz sąsiedztwa grafu, a~$A^{(2)} = A \cdot A$
— jej kwadrat (liczba ścieżek długości~2 między węzłami).
W~praktyce L3 obliczany jest jako element $(u,v)$ macierzy $A^3$.

\subsubsection{Katz}

Miara Katza uogólnia L3, sumując wkłady ścieżek wszystkich długości
z~wykładniczo malejącymi wagami~\cite{LibenNowell2007}:
%
\[
  \mathrm{Katz}(u, v) = \sum_{l=1}^{\infty} \beta^{l} \cdot |{\rm paths}_{uv}^{(l)}|
  = \left[(I - \beta A)^{-1} - I\right]_{uv},
\]
%
gdzie $\beta \in (0, 1)$ jest parametrem tłumienia, $A$ — macierzą
sąsiedztwa, a~$|{\rm paths}_{uv}^{(l)}|$ — liczbą ścieżek długości~$l$
między węzłami $u$ i~$v$. Mała wartość $\beta$ sprawia, że~dominujący
wkład pochodzi od~ścieżek najkrótszych. W~eksperymentach przyjęto
$\beta = 0{,}005$, co~zapewnia numeryczną stabilność obliczeń
przy~dużych grafach zbioru DLG-DG-23.

\subsubsection{Personalized PageRank}

\emph{Personalized PageRank} (PPR) realizuje ideę błądzenia losowego
z~restartem (\emph{random walk with restart})~\cite{LibenNowell2007}.
Dla~każdego węzła źródłowego $u$ oblicza się stacjonarny rozkład
prawdopodobieństwa błądzenia, które z~prawdopodobieństwem $\alpha$
powraca do~węzła $u$, a~z~prawdopodobieństwem $1 - \alpha$ przechodzi
do~losowego sąsiada. Wynik $\mathrm{PPR}(u, v)$ interpretuje się
jako prawdopodobieństwo odwiedzenia węzła $v$ podczas błądzenia
zainicjowanego w~$u$; wysoka wartość wskazuje na~bliskie sąsiedztwo
strukturalne. W~eksperymentach przyjęto $\alpha = 0{,}85$ (standardowa
wartość stosowana w~bibliotece NetworkX). Ze~względu na~koszt obliczenia
pełnej macierzy PPR dla~dużych grafów zastosowano obliczenia wsadowe
(\emph{batch computation}) z~wektorem startowym ograniczonym
do~węzłów zbioru treningowego.

% ============================================================
\subsection{Klasyczne metody uczenia maszynowego}
\label{sec:ml}
% ============================================================

Klasyczne metody uczenia maszynowego traktują problem wykrywania krawędzi
jako binarną klasyfikację, w~której wejściem jest wektor cech opisujący
parę węzłów $(u, v)$. Ze~względu na~anonimizację zbioru DLG-DG-23
(nazwy węzłów w~formacie UUID) jako cechy przyjęto wyłącznie atrybuty
topologiczne grafu — metryki oparte na~podobieństwie nazw (np.~Jaro-Winkler
stosowane w~\cite{Boinski2025}) są w~tym zbiorze nieaplikowalne.

\subsubsection{Wektor cech}

Dla~każdej pary kandydackiej $(u, v)$ wyznaczany jest wektor
11~cech topologicznych. Tabela~\ref{tab:cechy} przedstawia pełną
listę zastosowanych cech.

% TODO (WYMAGANE): Wstawić tabelę z opisem 11 cech:
%   Nr | Nazwa cechy | Opis | Uzasadnienie
%   Tabela musi być poprzedzona zdaniem „Tabela~X przedstawia...".
%   Uwzględnić: stopnie węzłów u i v, shortest_path, wyniki PA/L3/Katz/PPR,
%   przynależność do składowej spójnej, clustering coefficient itp.

Najwyższą istotność predykcyjną wykazuje cecha \emph{shortest\_path}
(długość najkrótszej ścieżki między $u$ a~$v$ w~grafie nieskierowanym),
co wskazuje, że~pary węzłów blisko połączonych w~istniejącej strukturze
grafu z~dużym prawdopodobieństwem są ze~sobą powiązane krawędzią
\emph{data\_flow} — nawet jeśli krawędź bezpośrednia nie~jest
obserwowana. Cecha ta~uzyskała ważność 52{,}6\% w~modelu Random Forest,
co~potwierdzono metodą permutacji na~zbiorze walidacyjnym.

\subsubsection{Random Forest}

Random Forest (RF) to~metoda uczenia zespołowego (\emph{ensemble
learning}) oparta na~uśrednianiu predykcji drzew decyzyjnych
wytrenowanych na~losowych podpróbkach danych i~cech
(\emph{bagging})~\cite{LibenNowell2007}. W~celu obsługi
niezrównoważonych klas zastosowano parametr
\texttt{class\_weight=\textquotedbl{}balanced\textquotedbl{}},
który skaluje wagi klas odwrotnie proporcjonalnie do~ich liczności.

\subsubsection{RUSBoost}

RUSBoost (\emph{Random Under-Sampling Boosting}) łączy algorytm
boostingu AdaBoost z~losowym podpróbkowaniem klasy większościowej
przed każdą iteracją. Metoda ta~jest szczególnie skuteczna przy silnym
niezbalansowaniu klas, typowym dla~problemu wykrywania krawędzi
w~rzadkich grafach. W~niniejszej pracy algorytm RUSBoost zaimplementowano
w~celu bezpośredniego porównania z~wynikami Boińskiego
et~al.~\cite{Boinski2025}, którzy uzyskali wartość F1~=~0{,}79
na~danych syntetycznych.

\subsubsection{LightGBM}

LightGBM to~implementacja gradientowego boostingu drzew decyzyjnych
(\emph{Gradient Boosted Decision Trees}) zoptymalizowana pod~kątem
wydajności i~skalowalności. W~odróżnieniu od~Random Forest, LightGBM
buduje kolejne drzewa w~sposób sekwencyjny, każde minimalizując błąd
resztkowy poprzednika. Parametr \texttt{scale\_pos\_weight} ustawiono
na~stosunek liczby próbek negatywnych do~pozytywnych, co~kompensuje
niezbalansowanie klas bez~konieczności podpróbowania.

% ============================================================
\subsection{Zanurzenia sieciowe — Node2Vec}
\label{sec:node2vec}
% ============================================================

\emph{Node2Vec} to~metoda uczenia reprezentacji węzłów grafu
(\emph{node embeddings}), opierająca się na~biased random walks
i~modelu Skip-gram~\cite{LibenNowell2007}. Każdemu węzłowi $v$
przypisywany jest wektor $\mathbf{z}_v \in \mathbb{R}^d$, zwany
zanurzeniem (\emph{embedding}), tak by~węzły o~podobnym kontekście
strukturalnym w~grafie były reprezentowane przez geometrycznie bliskie
wektory.

\subsubsection{Biased random walk}

Błądzenie losowe rozpoczyna się od~węzła $v$ i~na~każdym kroku
wybiera następny węzeł spośród sąsiadów bieżącego węzła
z~prawdopodobieństwem kontrolowanym przez dwa hiperparametry:

\begin{itemize}
  \item $p$ — parametr powrotu (\emph{return parameter}): kontroluje
        prawdopodobieństwo powrotu do~poprzednio odwiedzonego węzła;
        mała wartość $p$ sprzyja eksploatacji lokalnego sąsiedztwa,
  \item $q$ — parametr ekspansji (\emph{in-out parameter}): kontroluje
        skłonność do~oddalania się od~węzła startowego; wartość $q < 1$
        faworyzuje eksplorację dalszych obszarów grafu (strategia DFS),
        natomiast $q > 1$ — eksplorację lokalną (strategia BFS).
\end{itemize}

W~eksperymentach przyjęto $p = 1$ i~$q = 0{,}5$, co~sprzyja
eksploracji struktury globalnej grafu dwudzielnego — biased walk
z~$q < 1$ chętniej przemierza ścieżki prowadzące przez węzły obu
typów (\emph{job} i~\emph{table}), co~jest korzystne przy uczeniu
się reprezentacji par kandydackich.

\subsubsection{Generowanie embeddingu krawędzi}

Wykrywanie krawędzi wymaga przekształcenia pary zanurzeń węzłów
$(\mathbf{z}_u, \mathbf{z}_v)$ w~pojedynczy wektor reprezentujący
potencjalną krawędź $(u, v)$. Zastosowano iloczyn Hadamarda:
%
\[
  \mathbf{e}_{uv} = \mathbf{z}_u \odot \mathbf{z}_v,
\]
%
gdzie $\odot$ oznacza mnożenie element po~elemencie. Uzyskany
wektor $\mathbf{e}_{uv} \in \mathbb{R}^d$ podawany jest na~wejście
klasyfikatora — wielowarstwowej sieci neuronowej (ang.\
\emph{Multilayer Perceptron}, MLP) z~jedną warstwą ukrytą.

\subsubsection{Różnica względem grafowych sieci neuronowych}

Należy podkreślić, że~Node2Vec nie~jest grafową siecią neuronową
(ang.\ \emph{Graph Neural Network}, GNN). Metoda ta~uczy reprezentacji
węzłów wyłącznie na~podstawie sekwencji węzłów generowanych przez
błądzenie losowe (analogicznie do~metody Word2Vec na~sekwencjach
słów), bez~mechanizmu propagacji wiadomości (\emph{message passing})
charakterystycznego dla~GNN, takich jak GraphSAGE, R-GCN czy~HGT.
W~konsekwencji Node2Vec nie~uwzględnia atrybutów węzłów ani
typów krawędzi podczas uczenia zanurzeń. Rozszerzenie pracy
o~właściwe grafowe sieci neuronowe wskazano w~rozdziale~5 jako
kierunek dalszych badań.

% ============================================================

\chapter{Problem i~sposób jego rozwiązania}

Niniejszy rozdział precyzuje formalną definicję problemu broken lineage
i~jego ujęcie jako zadania klasyfikacji binarnej (sekcja~\ref{sec:sformulowanie}).
Sekcja~\ref{sec:architektura} opisuje architekturę zaproponowanego
potoku przetwarzania, a~sekcja~\ref{sec:dane} charakteryzuje zastosowany
zbiór danych. Metodykę symulacji brakujących krawędzi przedstawiono
w~sekcji~\ref{sec:scenariusze}, natomiast szczegóły implementacji
algorytmów omówiono w~sekcji~\ref{sec:algorytmy}.

% ============================================================
\section{Sformułowanie problemu}
\label{sec:sformulowanie}
% ============================================================

Graf lineage modeluje się jako graf skierowany $G^{*} = (V, E^{*})$,
gdzie $V$ oznacza zbiór obiektów bazodanowych (tabel, widoków, zadań ETL),
a~$E^{*}$ — zbiór wszystkich rzeczywistych relacji przepływu danych
(\emph{data\_flow}). W~praktyce część krawędzi jest niedostępna
z~powodu zjawiska broken lineage: tabele tymczasowe nie~są rejestrowane
w~standardowych słownikach bazy danych, a~funkcje definiowane przez
użytkownika (\emph{User-Defined Function}, UDF) ukrywają wewnętrzną
logikę transformacji. Obserwowany graf ma~zatem postać
$G_{\mathrm{obs}} = (V, E_{\mathrm{obs}})$, gdzie
$E_{\mathrm{obs}} \subsetneq E^{*}$. Brakujące krawędzie
$E^{*} \setminus E_{\mathrm{obs}}$ reprezentują zerwane zależności,
których odtworzenie jest celem niniejszej pracy.

Zadanie wykrywania brakujących krawędzi sprowadza się do~klasyfikacji
binarnej. Dla~każdej pary węzłów $(u, v) \in V \times V$ algorytm
wyznacza wartość $\hat{y}(u, v) \in [0, 1]$, interpretowaną jako
oszacowane prawdopodobieństwo istnienia krawędzi w~pełnym grafie
$G^{*}$. Para jest klasyfikowana jako krawędź istniejąca (klasa~1)
lub~nieistniejąca (klasa~0) po~przekroczeniu progu binaryzacji
wynoszącego $0{,}5$.

Ze~względu na~strukturę grafu DLG-DG-23 kandydatami do~wykrywania
są wyłącznie pary węzłów różnych typów: (\emph{job},~\emph{table})
lub~(\emph{table},~\emph{job}). Para (\emph{table},~\emph{table})
nigdy nie~jest połączona krawędzią \emph{data\_flow}, a~relacja
\emph{parent\_child} łączy wyłącznie węzły \emph{table} i~\emph{field}
— pary te nie~są przedmiotem wykrywania. Ograniczenie przestrzeni
kandydatów do~typologicznie dopuszczalnych par eliminuje trywialnie
fałszywe przykłady ujemne i~zapewnia spójność z~semantyką domeny.

% ============================================================
\section{Architektura rozwiązania}
\label{sec:architektura}
% ============================================================

Potok przetwarzania składa się z~czterech głównych komponentów:
modułu wczytywania danych (\emph{loader}), modułu podziału krawędzi
(\emph{splitter}), właściwych algorytmów wykrywania oraz modułu
ewaluacji.

% TODO (WYMAGANE): Wstawić rysunek blokowy potoku:
%   JSON → loader → NetworkX DiGraph → splitter → algorytmy → metryki.
%   Rysunek musi być wprowadzony zdaniem przed jego wystąpieniem,
%   np. „Rysunek~X przedstawia schemat blokowy zaproponowanego potoku."

\subsection{Moduł wczytywania danych}

Dane wejściowe dla~każdego grafu DLG przechowywane są w~dwóch plikach
JSON: \texttt{DLGi-node.json} zawiera opisy węzłów wraz z~atrybutem
\texttt{asset\_type}, natomiast \texttt{DLGi-edge.json} — opisy
krawędzi z~atrybutami \texttt{relation\_type} i~\texttt{relation\_id}.
Moduł \texttt{loader.py} wczytuje oba pliki dla~każdego grafu i~buduje
obiekt \texttt{nx.DiGraph} z~biblioteki NetworkX, a~następnie zwraca
słownik wszystkich 18~grafów zbioru DLG-DG-23. Sam graf przechowuje
atrybut \texttt{name} (np.\ \texttt{DLG1}).

\subsection{Moduł podziału krawędzi}

Moduł \texttt{splitter.py} wyodrębnia ze~struktury grafu krawędzie
typu \emph{data\_flow} i~dzieli je~na~trzy rozłączne podzbiory
w~proporcji 60/20/20 (treningowy, walidacyjny, testowy). Próbki
pozytywne stanowią krawędzie usunięte ze~zbioru treningowego
i~przeznaczone do~ewaluacji; próbki negatywne to~pary węzłów
typologicznie dopuszczalnych (wyłącznie \emph{job}--\emph{table}),
które nigdy nie~były połączone krawędzią. Stosunek próbek negatywnych
do~pozytywnych wynosi 1:1, co~zapewnia zbalansowany zbiór ewaluacyjny
i~stabilne wartości metryk porównawczych.

Moduł \texttt{scenario\_splitter.py} realizuje analogiczny podział,
lecz zamiast losowego usuwania krawędzi stosuje trzy scenariusze
symulacji broken lineage opisane w~sekcji~\ref{sec:scenariusze}.

\subsection{Moduł ewaluacji}

Moduł \texttt{metrics.py} wyznacza standardowe metryki klasyfikacji
binarnej: precyzję, czułość, miarę F1, AUC-ROC i~AUC-PR. Próg
binaryzacji wynosi $0{,}5$ dla~wszystkich metod. Wyniki są
uśredniane po~wszystkich 18~grafach zbioru DLG-DG-23 i~agregowane
według scenariusza oraz algorytmu, co~umożliwia bezpośrednie
porównanie między metodami.

% ============================================================
\section{Zbiór danych DLG-DG-23}
\label{sec:dane}
% ============================================================

Zastosowany zbiór danych pochodzi z~pracy Chen et~al.~\cite{Chen2024}
i~zawiera 18~grafów lineage pobranych ze~środowiska produkcyjnego
Huawei Cloud. Dane zostały zanonimizowane — nazwy tabel i~atrybutów
zastąpiono identyfikatorami w~formacie UUID, a~schematy bazy danych
nie~zostały udostępnione. Zbiór opublikowano w~2024~roku jako pierwszy
otwarty zestaw rzeczywistych grafów lineage; szczegółową charakterystykę
poszczególnych grafów zamieszczono w~tabeli~\ref{tab:dlg-stats} w~sekcji~\ref{sec:dane-opis}.

Zbiór DLG-DG-23 wybrano ze~względu na~jego unikatowy charakter:
jest to~jedyny publicznie dostępny zestaw rzeczywistych grafów lineage,
co~umożliwia ewaluację algorytmów na~danych o~autentycznej topologii
przemysłowej, bez~ryzyka artefaktów wprowadzonych przez dane
syntetyczne. Anonimizacja nazw węzłów ma~bezpośredni wpływ na~dobór
algorytmów: metryki podobieństwa oparte na~nazwach obiektów
(np.~odległość Jaro--Winkler stosowana przez Boińskiego
et~al.~\cite{Boinski2025}) są nieaplikowalne. W~niniejszej pracy
wektor cech ograniczono zatem do~atrybutów topologicznych grafu,
co~jednocześnie stanowi element badań nad~pytaniem, czy sama struktura
grafu lineage jest wystarczająca do~wykrywania zerwanych zależności.

\subsection{Charakterystyka grafów}
\label{sec:dane-opis}

Tabela~\ref{tab:dlg-stats} przedstawia statystyki poszczególnych
grafów zbioru DLG-DG-23.

\begin{table}[ht]
\centering
\caption{Charakterystyka grafów zbioru DLG-DG-23.}
\label{tab:dlg-stats}
\begin{threeparttable}
\begin{tabular}{llrl}
\toprule
\textbf{ID} & \textbf{Węzły} & \textbf{Krawędzie DATA\_FLOW} & \textbf{Skala} \\
\midrule
DLG1  &    298 &    24 & mały   \\
DLG2  &    464 &    37 & mały   \\
DLG3  &    603 &    57 & mały   \\
DLG4  &    415 &    32 & mały   \\
DLG5  & 1\,299 &   175 & średni \\
DLG6  & 13\,840 & 1\,443 & duży  \\
DLG7  &    574 &    43 & mały   \\
DLG8  &    546 &    34 & mały   \\
DLG9  &    756 &    62 & mały   \\
DLG10 & 5\,378 &   497 & duży   \\
DLG11 & 2\,024 &   121 & średni \\
DLG12 & 5\,645 &   469 & duży   \\
DLG13 & 2\,157 &   229 & średni \\
DLG14 &    453 &    24 & mały   \\
DLG15 &    558 &    41 & mały   \\
DLG16 & 5\,574 &   745 & duży   \\
DLG17 &    651 &    38 & mały   \\
DLG18 & 17\,085 & 1\,694 & duży \\
\bottomrule
\end{tabular}
\begin{tablenotes}
\small
\item Skala: mały $<$ 1\,000 węzłów; średni 1\,000–5\,000; duży $>$ 5\,000.
\item Liczba krawędzi DATA\_FLOW podana dla sumy zbiorów treningowego,
      walidacyjnego i~testowego (podział 60/20/20).
\end{tablenotes}
\end{threeparttable}
\end{table}

% ============================================================
\section{Metodyka symulacji broken lineage}
\label{sec:scenariusze}
% ============================================================

Ponieważ zbiór DLG-DG-23 zawiera w~pełni zaobserwowane grafy lineage,
zjawisko broken lineage symulowane jest poprzez kontrolowane usuwanie
krawędzi. Podejście to jest standardowe w~literaturze wykrywania
krawędzi~\cite{LibenNowell2007}: dysponując kompletnym grafem
$G^{*}$, usuwa się podzbiór krawędzi, tworząc graf obserwowany
$G_{\mathrm{obs}}$, a~następnie sprawdza, czy algorytm potrafi
odtworzyć usunięte krawędzie. W~niniejszej pracy przyjęto trzy
scenariusze symulacji, odzwierciedlające typowe przyczyny powstawania
broken lineage.

\textbf{Scenariusz A} — losowe usunięcie krawędzi — realizuje
najprostszą formę symulacji: 20\% krawędzi \emph{data\_flow} jest
usuwanych ze~zbioru treningowego w~sposób losowy, bez~względu na~typ
węzłów końcowych. Scenariusz ten modeluje sytuację, w~której~broken
lineage wynika z~przyczyn losowych (np.~niedokompletowania logów
systemowych) i~nie~jest skoncentrowany na~konkretnym typie obiektu.

\textbf{Scenariusz B} — usunięcie krawędzi tabeli pośredniej —
symuluje sytuację, w~której tabela pośrednia (staging) jest tabelą
tymczasową i~nie~zostaje zarejestrowana w~metadanych. Usuwane są
wszystkie krawędzie prowadzące do~wybranego węzła \emph{table}
o~wysokim stopniu (tabela pośrednia), co~powoduje odcięcie
całego podgrafu zależności przepływającego przez ten węzeł.

\textbf{Scenariusz C} — usunięcie krawędzi do~tabeli docelowej —
symuluje ukrycie zależności na~ostatnim etapie potoku ETL, kiedy
krawędź do~tabeli docelowej (\emph{data mart}) nie~zostaje utrwalona.
Usuwane są krawędzie wchodzące do~wybranego węzła \emph{table}
o~zerowym lub~niskim stopniu wyjściowym, co~odpowiada tabelom
końcowym w~hierarchii przepływu danych.

% TODO (WYMAGANE): Wstawić tabelę zbiorczą scenariuszy:
%   Scenariusz | Co usuwamy | Co symuluje | Typ węzłów dotkniętych
%   Tabela musi być poprzedzona zdaniem „Tabela X przedstawia...".

% TODO (OPCJONALNIE): Wstawić poglądowe rysunki małego grafu
%   przed i~po usunięciu krawędzi dla~każdego scenariusza.
%   Rysunki muszą być wprowadzone zdaniem w~tekście.

% ============================================================
\section{Zaimplementowane algorytmy}
\label{sec:algorytmy}
% ============================================================

W~pracy zaimplementowano trzy grupy algorytmów wykrywania
brakujących krawędzi. Pierwszą grupę stanowią heurystyki grafowe
— metody niewymagające uczenia, opierające się wyłącznie na~strukturze
grafu. Drugą grupę tworzą klasyczne metody uczenia maszynowego,
budujące klasyfikator na~podstawie wektora cech topologicznych
pary węzłów. Trzecia grupa obejmuje metody zanurzeń sieciowych
(\emph{network embeddings}), reprezentowane przez algorytm Node2Vec.

% ============================================================
\subsection{Heurystyki grafowe}
% ============================================================

Heurystyki grafowe stanowią najprostszą klasę metod wykrywania krawędzi
(\emph{link prediction}). Wynik dla~pary węzłów $(u, v)$ wyznaczany jest
wyłącznie na~podstawie lokalnej lub~globalnej struktury grafu, bez~etapu
trenowania modelu. Metody te charakteryzują się niską złożonością
obliczeniową i~pełną interpretowalnością, co~czyni je~naturalnym punktem
odniesienia (ang.\ \emph{baseline}) przy~porównywaniu bardziej
zaawansowanych podejść~\cite{LibenNowell2007}.

Wstępna analiza struktury grafów DLG-DG-23 wykazała, że~miary oparte
na~wspólnych sąsiadach — Common Neighbors, indeks Jaccarda
i~miara Adamic-Adar — są nieskuteczne w~grafach dwudzielnych.
W~strukturze bipartite węzły tego samego typu (\emph{job}--\emph{job}
lub~\emph{table}--\emph{table}) nigdy nie~posiadają wspólnych sąsiadów
bezpośrednich, ponieważ krawędzie \emph{data\_flow} łączą wyłącznie
węzły różnych typów. W~konsekwencji wymienione miary zwracają wynik
zerowy dla~każdej typologicznie dopuszczalnej pary kandydackiej,
co~dyskwalifikuje je~jako metody wykrywania w~omawianym scenariuszu.
W~implementacji uwzględniono wyłącznie heurystyki działające poprawnie
na~grafach dwudzielnych: Preferential Attachment, L3, Katz
oraz~Personalized PageRank.

\subsubsection{Preferential Attachment}

Miara \emph{Preferential Attachment} (PA) opiera się na~obserwacji,
że~w~sieciach rzeczywistych nowe połączenia powstają z~prawdopodobieństwem
proporcjonalnym do~stopnia istniejących węzłów — węzły o~wielu sąsiadach
z~większym prawdopodobieństwem pozyskują kolejne
połączenia~\cite{Barabasi1999}. Jako heurystyka wykrywania krawędzi PA
została zastosowana przez Liben-Nowella i~Kleinberga~\cite{LibenNowell2007}.

Dla~pary węzłów $(u, v)$ wynik PA definiuje się jako:
%
\[
  \mathrm{PA}(u, v) = |N(u)| \cdot |N(v)|,
\]
%
gdzie $N(u)$ oraz $N(v)$ oznaczają zbiory sąsiadów odpowiednio węzła $u$
i~węzła $v$, a~$|\cdot|$ oznacza rozmiar zbioru, tj.\ stopień węzła w~grafie.

W~kontekście grafów lineage wysoki stopień węzła \emph{job} oznacza,
że~dane zadanie czyta z~wielu tabel źródłowych lub~zapisuje do~wielu
tabel docelowych; wysoki stopień węzła \emph{table} wskazuje na~tabelę
centralną — często pośrednią warstwę integracyjną w~potoku ETL.
Para $(u, v)$, w~której oba węzły mają duże stopnie, otrzymuje wysoki
wynik PA, co~interpretuje się jako sygnał potencjalnej zależności
\emph{data\_flow}. Miara PA nie~wymaga istnienia wspólnych sąsiadów,
co~stanowi jej kluczową przewagę w~grafach dwudzielnych.

\subsubsection{L3}

Miara L3 (ang.\ \emph{path of length 3}) liczy ścieżki długości~3
między węzłami $u$ i~$v$~\cite{LibenNowell2007}. W~grafie dwudzielnym
ścieżka długości~3 między węzłem \emph{job} a~węzłem \emph{table}
przebiega przez pośredni węzeł \emph{table} i~pośredni węzeł \emph{job}
— co odpowiada wzorcowi: zadanie~$u$ czyta z~tabeli~$t_1$, tabela~$t_1$
jest zapisywana przez zadanie~$j_1$, zadanie~$j_1$ czyta z~tabeli~$v$.
Wysoka wartość L3 sygnalizuje, że~para $(u, v)$ jest gęsto połączona
przez pośrednie obiekty, co~w~kontekście broken lineage może wskazywać
na~zależność ukrytą przez tabelę tymczasową.

Wynik L3 dla~pary $(u, v)$ definiuje się jako:
%
\[
  \mathrm{L3}(u, v) = \sum_{w \in V} A_{uw} \cdot A_{wv}^{(2)},
\]
%
gdzie $A$ oznacza macierz sąsiedztwa grafu, a~$A^{(2)} = A \cdot A$
— jej kwadrat (liczba ścieżek długości~2 między węzłami).
W~praktyce L3 obliczany jest jako element $(u,v)$ macierzy $A^3$.

\subsubsection{Katz}

Miara Katza uogólnia L3, sumując wkłady ścieżek wszystkich długości
z~wykładniczo malejącymi wagami~\cite{LibenNowell2007}:
%
\[
  \mathrm{Katz}(u, v) = \sum_{l=1}^{\infty} \beta^{l} \cdot |{\rm paths}_{uv}^{(l)}|
  = \left[(I - \beta A)^{-1} - I\right]_{uv},
\]
%
gdzie $\beta \in (0, 1)$ jest parametrem tłumienia, $A$ — macierzą
sąsiedztwa, a~$|{\rm paths}_{uv}^{(l)}|$ — liczbą ścieżek długości~$l$
między węzłami $u$ i~$v$. Mała wartość $\beta$ sprawia, że~dominujący
wkład pochodzi od~ścieżek najkrótszych. W~eksperymentach przyjęto
$\beta = 0{,}005$, co~zapewnia numeryczną stabilność obliczeń
przy~dużych grafach zbioru DLG-DG-23.

\subsubsection{Personalized PageRank}

\emph{Personalized PageRank} (PPR) realizuje ideę błądzenia losowego
z~restartem (\emph{random walk with restart})~\cite{LibenNowell2007}.
Dla~każdego węzła źródłowego $u$ oblicza się stacjonarny rozkład
prawdopodobieństwa błądzenia, które z~prawdopodobieństwem $\alpha$
powraca do~węzła $u$, a~z~prawdopodobieństwem $1 - \alpha$ przechodzi
do~losowego sąsiada. Wynik $\mathrm{PPR}(u, v)$ interpretuje się
jako prawdopodobieństwo odwiedzenia węzła $v$ podczas błądzenia
zainicjowanego w~$u$; wysoka wartość wskazuje na~bliskie sąsiedztwo
strukturalne. W~eksperymentach przyjęto $\alpha = 0{,}85$ (standardowa
wartość stosowana w~bibliotece NetworkX). Ze~względu na~koszt obliczenia
pełnej macierzy PPR dla~dużych grafów zastosowano obliczenia wsadowe
(\emph{batch computation}) z~wektorem startowym ograniczonym
do~węzłów zbioru treningowego.

% ============================================================
\subsection{Klasyczne metody uczenia maszynowego}
\label{sec:ml}
% ============================================================

Klasyczne metody uczenia maszynowego traktują problem wykrywania krawędzi
jako binarną klasyfikację, w~której wejściem jest wektor cech opisujący
parę węzłów $(u, v)$. Ze~względu na~anonimizację zbioru DLG-DG-23
(nazwy węzłów w~formacie UUID) jako cechy przyjęto wyłącznie atrybuty
topologiczne grafu — metryki oparte na~podobieństwie nazw (np.~Jaro-Winkler
stosowane w~\cite{Boinski2025}) są w~tym zbiorze nieaplikowalne.

\subsubsection{Wektor cech}

Dla~każdej pary kandydackiej $(u, v)$ wyznaczany jest wektor
11~cech topologicznych. Tabela~\ref{tab:cechy} przedstawia pełną
listę zastosowanych cech.

% TODO (WYMAGANE): Wstawić tabelę z opisem 11 cech:
%   Nr | Nazwa cechy | Opis | Uzasadnienie
%   Tabela musi być poprzedzona zdaniem „Tabela~X przedstawia...".
%   Uwzględnić: stopnie węzłów u i v, shortest_path, wyniki PA/L3/Katz/PPR,
%   przynależność do składowej spójnej, clustering coefficient itp.

Najwyższą istotność predykcyjną wykazuje cecha \emph{shortest\_path}
(długość najkrótszej ścieżki między $u$ a~$v$ w~grafie nieskierowanym),
co wskazuje, że~pary węzłów blisko połączonych w~istniejącej strukturze
grafu z~dużym prawdopodobieństwem są ze~sobą powiązane krawędzią
\emph{data\_flow} — nawet jeśli krawędź bezpośrednia nie~jest
obserwowana. Cecha ta~uzyskała ważność 52{,}6\% w~modelu Random Forest,
co~potwierdzono metodą permutacji na~zbiorze walidacyjnym.

\subsubsection{Random Forest}

Random Forest (RF) to~metoda uczenia zespołowego (\emph{ensemble
learning}) oparta na~uśrednianiu predykcji drzew decyzyjnych
wytrenowanych na~losowych podpróbkach danych i~cech
(\emph{bagging})~\cite{LibenNowell2007}. W~celu obsługi
niezrównoważonych klas zastosowano parametr
\texttt{class\_weight=\textquotedbl{}balanced\textquotedbl{}},
który skaluje wagi klas odwrotnie proporcjonalnie do~ich liczności.

\subsubsection{RUSBoost}

RUSBoost (\emph{Random Under-Sampling Boosting}) łączy algorytm
boostingu AdaBoost z~losowym podpróbkowaniem klasy większościowej
przed każdą iteracją. Metoda ta~jest szczególnie skuteczna przy silnym
niezbalansowaniu klas, typowym dla~problemu wykrywania krawędzi
w~rzadkich grafach. W~niniejszej pracy algorytm RUSBoost zaimplementowano
w~celu bezpośredniego porównania z~wynikami Boińskiego
et~al.~\cite{Boinski2025}, którzy uzyskali wartość F1~=~0{,}79
na~danych syntetycznych.

\subsubsection{LightGBM}

LightGBM to~implementacja gradientowego boostingu drzew decyzyjnych
(\emph{Gradient Boosted Decision Trees}) zoptymalizowana pod~kątem
wydajności i~skalowalności. W~odróżnieniu od~Random Forest, LightGBM
buduje kolejne drzewa w~sposób sekwencyjny, każde minimalizując błąd
resztkowy poprzednika. Parametr \texttt{scale\_pos\_weight} ustawiono
na~stosunek liczby próbek negatywnych do~pozytywnych, co~kompensuje
niezbalansowanie klas bez~konieczności podpróbowania.

% ============================================================
\subsection{Zanurzenia sieciowe — Node2Vec}
\label{sec:node2vec}
% ============================================================

\emph{Node2Vec} to~metoda uczenia reprezentacji węzłów grafu
(\emph{node embeddings}), opierająca się na~biased random walks
i~modelu Skip-gram~\cite{LibenNowell2007}. Każdemu węzłowi $v$
przypisywany jest wektor $\mathbf{z}_v \in \mathbb{R}^d$, zwany
zanurzeniem (\emph{embedding}), tak by~węzły o~podobnym kontekście
strukturalnym w~grafie były reprezentowane przez geometrycznie bliskie
wektory.

\subsubsection{Biased random walk}

Błądzenie losowe rozpoczyna się od~węzła $v$ i~na~każdym kroku
wybiera następny węzeł spośród sąsiadów bieżącego węzła
z~prawdopodobieństwem kontrolowanym przez dwa hiperparametry:

\begin{itemize}
  \item $p$ — parametr powrotu (\emph{return parameter}): kontroluje
        prawdopodobieństwo powrotu do~poprzednio odwiedzonego węzła;
        mała wartość $p$ sprzyja eksploatacji lokalnego sąsiedztwa,
  \item $q$ — parametr ekspansji (\emph{in-out parameter}): kontroluje
        skłonność do~oddalania się od~węzła startowego; wartość $q < 1$
        faworyzuje eksplorację dalszych obszarów grafu (strategia DFS),
        natomiast $q > 1$ — eksplorację lokalną (strategia BFS).
\end{itemize}

W~eksperymentach przyjęto $p = 1$ i~$q = 0{,}5$, co~sprzyja
eksploracji struktury globalnej grafu dwudzielnego — biased walk
z~$q < 1$ chętniej przemierza ścieżki prowadzące przez węzły obu
typów (\emph{job} i~\emph{table}), co~jest korzystne przy uczeniu
się reprezentacji par kandydackich.

\subsubsection{Generowanie embeddingu krawędzi}

Wykrywanie krawędzi wymaga przekształcenia pary zanurzeń węzłów
$(\mathbf{z}_u, \mathbf{z}_v)$ w~pojedynczy wektor reprezentujący
potencjalną krawędź $(u, v)$. Zastosowano iloczyn Hadamarda:
%
\[
  \mathbf{e}_{uv} = \mathbf{z}_u \odot \mathbf{z}_v,
\]
%
gdzie $\odot$ oznacza mnożenie element po~elemencie. Uzyskany
wektor $\mathbf{e}_{uv} \in \mathbb{R}^d$ podawany jest na~wejście
klasyfikatora — wielowarstwowej sieci neuronowej (ang.\
\emph{Multilayer Perceptron}, MLP) z~jedną warstwą ukrytą.

\subsubsection{Różnica względem grafowych sieci neuronowych}

Należy podkreślić, że~Node2Vec nie~jest grafową siecią neuronową
(ang.\ \emph{Graph Neural Network}, GNN). Metoda ta~uczy reprezentacji
węzłów wyłącznie na~podstawie sekwencji węzłów generowanych przez
błądzenie losowe (analogicznie do~metody Word2Vec na~sekwencjach
słów), bez~mechanizmu propagacji wiadomości (\emph{message passing})
charakterystycznego dla~GNN, takich jak GraphSAGE, R-GCN czy~HGT.
W~konsekwencji Node2Vec nie~uwzględnia atrybutów węzłów ani
typów krawędzi podczas uczenia zanurzeń. Rozszerzenie pracy
o~właściwe grafowe sieci neuronowe wskazano w~rozdziale~5 jako
kierunek dalszych badań.

% ============================================================

\chapter{Problem i~sposób jego rozwiązania}

Niniejszy rozdział precyzuje formalną definicję problemu broken lineage
i~jego ujęcie jako zadania klasyfikacji binarnej (sekcja~\ref{sec:sformulowanie}).
Sekcja~\ref{sec:architektura} opisuje architekturę zaproponowanego
potoku przetwarzania, a~sekcja~\ref{sec:dane} charakteryzuje zastosowany
zbiór danych. Metodykę symulacji brakujących krawędzi przedstawiono
w~sekcji~\ref{sec:scenariusze}, natomiast szczegóły implementacji
algorytmów omówiono w~sekcji~\ref{sec:algorytmy}.

% ============================================================
\section{Sformułowanie problemu}
\label{sec:sformulowanie}
% ============================================================

Graf lineage modeluje się jako graf skierowany $G^{*} = (V, E^{*})$,
gdzie $V$ oznacza zbiór obiektów bazodanowych (tabel, widoków, zadań ETL),
a~$E^{*}$ — zbiór wszystkich rzeczywistych relacji przepływu danych
(\emph{data\_flow}). W~praktyce część krawędzi jest niedostępna
z~powodu zjawiska broken lineage: tabele tymczasowe nie~są rejestrowane
w~standardowych słownikach bazy danych, a~funkcje definiowane przez
użytkownika (\emph{User-Defined Function}, UDF) ukrywają wewnętrzną
logikę transformacji. Obserwowany graf ma~zatem postać
$G_{\mathrm{obs}} = (V, E_{\mathrm{obs}})$, gdzie
$E_{\mathrm{obs}} \subsetneq E^{*}$. Brakujące krawędzie
$E^{*} \setminus E_{\mathrm{obs}}$ reprezentują zerwane zależności,
których odtworzenie jest celem niniejszej pracy.

Zadanie wykrywania brakujących krawędzi sprowadza się do~klasyfikacji
binarnej. Dla~każdej pary węzłów $(u, v) \in V \times V$ algorytm
wyznacza wartość $\hat{y}(u, v) \in [0, 1]$, interpretowaną jako
oszacowane prawdopodobieństwo istnienia krawędzi w~pełnym grafie
$G^{*}$. Para jest klasyfikowana jako krawędź istniejąca (klasa~1)
lub~nieistniejąca (klasa~0) po~przekroczeniu progu binaryzacji
wynoszącego $0{,}5$.

Ze~względu na~strukturę grafu DLG-DG-23 kandydatami do~wykrywania
są wyłącznie pary węzłów różnych typów: (\emph{job},~\emph{table})
lub~(\emph{table},~\emph{job}). Para (\emph{table},~\emph{table})
nigdy nie~jest połączona krawędzią \emph{data\_flow}, a~relacja
\emph{parent\_child} łączy wyłącznie węzły \emph{table} i~\emph{field}
— pary te nie~są przedmiotem wykrywania. Ograniczenie przestrzeni
kandydatów do~typologicznie dopuszczalnych par eliminuje trywialnie
fałszywe przykłady ujemne i~zapewnia spójność z~semantyką domeny.

% ============================================================
\section{Architektura rozwiązania}
\label{sec:architektura}
% ============================================================

Potok przetwarzania składa się z~czterech głównych komponentów:
modułu wczytywania danych (\emph{loader}), modułu podziału krawędzi
(\emph{splitter}), właściwych algorytmów wykrywania oraz modułu
ewaluacji.

% TODO (WYMAGANE): Wstawić rysunek blokowy potoku:
%   JSON → loader → NetworkX DiGraph → splitter → algorytmy → metryki.
%   Rysunek musi być wprowadzony zdaniem przed jego wystąpieniem,
%   np. „Rysunek~X przedstawia schemat blokowy zaproponowanego potoku."

\subsection{Moduł wczytywania danych}

Dane wejściowe dla~każdego grafu DLG przechowywane są w~dwóch plikach
JSON: \texttt{DLGi-node.json} zawiera opisy węzłów wraz z~atrybutem
\texttt{asset\_type}, natomiast \texttt{DLGi-edge.json} — opisy
krawędzi z~atrybutami \texttt{relation\_type} i~\texttt{relation\_id}.
Moduł \texttt{loader.py} wczytuje oba pliki dla~każdego grafu i~buduje
obiekt \texttt{nx.DiGraph} z~biblioteki NetworkX, a~następnie zwraca
słownik wszystkich 18~grafów zbioru DLG-DG-23. Sam graf przechowuje
atrybut \texttt{name} (np.\ \texttt{DLG1}).

\subsection{Moduł podziału krawędzi}

Moduł \texttt{splitter.py} wyodrębnia ze~struktury grafu krawędzie
typu \emph{data\_flow} i~dzieli je~na~trzy rozłączne podzbiory
w~proporcji 60/20/20 (treningowy, walidacyjny, testowy). Próbki
pozytywne stanowią krawędzie usunięte ze~zbioru treningowego
i~przeznaczone do~ewaluacji; próbki negatywne to~pary węzłów
typologicznie dopuszczalnych (wyłącznie \emph{job}--\emph{table}),
które nigdy nie~były połączone krawędzią. Stosunek próbek negatywnych
do~pozytywnych wynosi 1:1, co~zapewnia zbalansowany zbiór ewaluacyjny
i~stabilne wartości metryk porównawczych.

Moduł \texttt{scenario\_splitter.py} realizuje analogiczny podział,
lecz zamiast losowego usuwania krawędzi stosuje trzy scenariusze
symulacji broken lineage opisane w~sekcji~\ref{sec:scenariusze}.

\subsection{Moduł ewaluacji}

Moduł \texttt{metrics.py} wyznacza standardowe metryki klasyfikacji
binarnej: precyzję, czułość, miarę F1, AUC-ROC i~AUC-PR. Próg
binaryzacji wynosi $0{,}5$ dla~wszystkich metod. Wyniki są
uśredniane po~wszystkich 18~grafach zbioru DLG-DG-23 i~agregowane
według scenariusza oraz algorytmu, co~umożliwia bezpośrednie
porównanie między metodami.

% ============================================================
\section{Zbiór danych DLG-DG-23}
\label{sec:dane}
% ============================================================

Zastosowany zbiór danych pochodzi z~pracy Chen et~al.~\cite{Chen2024}
i~zawiera 18~grafów lineage pobranych ze~środowiska produkcyjnego
Huawei Cloud. Dane zostały zanonimizowane — nazwy tabel i~atrybutów
zastąpiono identyfikatorami w~formacie UUID, a~schematy bazy danych
nie~zostały udostępnione. Zbiór opublikowano w~2024~roku jako pierwszy
otwarty zestaw rzeczywistych grafów lineage; szczegółową charakterystykę
poszczególnych grafów zamieszczono w~tabeli~\ref{tab:dlg-stats} w~sekcji~\ref{sec:dane-opis}.

Zbiór DLG-DG-23 wybrano ze~względu na~jego unikatowy charakter:
jest to~jedyny publicznie dostępny zestaw rzeczywistych grafów lineage,
co~umożliwia ewaluację algorytmów na~danych o~autentycznej topologii
przemysłowej, bez~ryzyka artefaktów wprowadzonych przez dane
syntetyczne. Anonimizacja nazw węzłów ma~bezpośredni wpływ na~dobór
algorytmów: metryki podobieństwa oparte na~nazwach obiektów
(np.~odległość Jaro--Winkler stosowana przez Boińskiego
et~al.~\cite{Boinski2025}) są nieaplikowalne. W~niniejszej pracy
wektor cech ograniczono zatem do~atrybutów topologicznych grafu,
co~jednocześnie stanowi element badań nad~pytaniem, czy sama struktura
grafu lineage jest wystarczająca do~wykrywania zerwanych zależności.

\subsection{Charakterystyka grafów}
\label{sec:dane-opis}

Tabela~\ref{tab:dlg-stats} przedstawia statystyki poszczególnych
grafów zbioru DLG-DG-23.

\begin{table}[ht]
\centering
\caption{Charakterystyka grafów zbioru DLG-DG-23.}
\label{tab:dlg-stats}
\begin{threeparttable}
\begin{tabular}{llrl}
\toprule
\textbf{ID} & \textbf{Węzły} & \textbf{Krawędzie DATA\_FLOW} & \textbf{Skala} \\
\midrule
DLG1  &    298 &    24 & mały   \\
DLG2  &    464 &    37 & mały   \\
DLG3  &    603 &    57 & mały   \\
DLG4  &    415 &    32 & mały   \\
DLG5  & 1\,299 &   175 & średni \\
DLG6  & 13\,840 & 1\,443 & duży  \\
DLG7  &    574 &    43 & mały   \\
DLG8  &    546 &    34 & mały   \\
DLG9  &    756 &    62 & mały   \\
DLG10 & 5\,378 &   497 & duży   \\
DLG11 & 2\,024 &   121 & średni \\
DLG12 & 5\,645 &   469 & duży   \\
DLG13 & 2\,157 &   229 & średni \\
DLG14 &    453 &    24 & mały   \\
DLG15 &    558 &    41 & mały   \\
DLG16 & 5\,574 &   745 & duży   \\
DLG17 &    651 &    38 & mały   \\
DLG18 & 17\,085 & 1\,694 & duży \\
\bottomrule
\end{tabular}
\begin{tablenotes}
\small
\item Skala: mały $<$ 1\,000 węzłów; średni 1\,000–5\,000; duży $>$ 5\,000.
\item Liczba krawędzi DATA\_FLOW podana dla sumy zbiorów treningowego,
      walidacyjnego i~testowego (podział 60/20/20).
\end{tablenotes}
\end{threeparttable}
\end{table}

% ============================================================
\section{Metodyka symulacji broken lineage}
\label{sec:scenariusze}
% ============================================================

Ponieważ zbiór DLG-DG-23 zawiera w~pełni zaobserwowane grafy lineage,
zjawisko broken lineage symulowane jest poprzez kontrolowane usuwanie
krawędzi. Podejście to jest standardowe w~literaturze wykrywania
krawędzi~\cite{LibenNowell2007}: dysponując kompletnym grafem
$G^{*}$, usuwa się podzbiór krawędzi, tworząc graf obserwowany
$G_{\mathrm{obs}}$, a~następnie sprawdza, czy algorytm potrafi
odtworzyć usunięte krawędzie. W~niniejszej pracy przyjęto trzy
scenariusze symulacji, odzwierciedlające typowe przyczyny powstawania
broken lineage.

\textbf{Scenariusz A} — losowe usunięcie krawędzi — realizuje
najprostszą formę symulacji: 20\% krawędzi \emph{data\_flow} jest
usuwanych ze~zbioru treningowego w~sposób losowy, bez~względu na~typ
węzłów końcowych. Scenariusz ten modeluje sytuację, w~której~broken
lineage wynika z~przyczyn losowych (np.~niedokompletowania logów
systemowych) i~nie~jest skoncentrowany na~konkretnym typie obiektu.

\textbf{Scenariusz B} — usunięcie krawędzi tabeli pośredniej —
symuluje sytuację, w~której tabela pośrednia (staging) jest tabelą
tymczasową i~nie~zostaje zarejestrowana w~metadanych. Usuwane są
wszystkie krawędzie prowadzące do~wybranego węzła \emph{table}
o~wysokim stopniu (tabela pośrednia), co~powoduje odcięcie
całego podgrafu zależności przepływającego przez ten węzeł.

\textbf{Scenariusz C} — usunięcie krawędzi do~tabeli docelowej —
symuluje ukrycie zależności na~ostatnim etapie potoku ETL, kiedy
krawędź do~tabeli docelowej (\emph{data mart}) nie~zostaje utrwalona.
Usuwane są krawędzie wchodzące do~wybranego węzła \emph{table}
o~zerowym lub~niskim stopniu wyjściowym, co~odpowiada tabelom
końcowym w~hierarchii przepływu danych.

% TODO (WYMAGANE): Wstawić tabelę zbiorczą scenariuszy:
%   Scenariusz | Co usuwamy | Co symuluje | Typ węzłów dotkniętych
%   Tabela musi być poprzedzona zdaniem „Tabela X przedstawia...".

% TODO (OPCJONALNIE): Wstawić poglądowe rysunki małego grafu
%   przed i~po usunięciu krawędzi dla~każdego scenariusza.
%   Rysunki muszą być wprowadzone zdaniem w~tekście.

% ============================================================
\section{Zaimplementowane algorytmy}
\label{sec:algorytmy}
% ============================================================

W~pracy zaimplementowano trzy grupy algorytmów wykrywania
brakujących krawędzi. Pierwszą grupę stanowią heurystyki grafowe
— metody niewymagające uczenia, opierające się wyłącznie na~strukturze
grafu. Drugą grupę tworzą klasyczne metody uczenia maszynowego,
budujące klasyfikator na~podstawie wektora cech topologicznych
pary węzłów. Trzecia grupa obejmuje metody zanurzeń sieciowych
(\emph{network embeddings}), reprezentowane przez algorytm Node2Vec.

% ============================================================
\subsection{Heurystyki grafowe}
% ============================================================

Heurystyki grafowe stanowią najprostszą klasę metod wykrywania krawędzi
(\emph{link prediction}). Wynik dla~pary węzłów $(u, v)$ wyznaczany jest
wyłącznie na~podstawie lokalnej lub~globalnej struktury grafu, bez~etapu
trenowania modelu. Metody te charakteryzują się niską złożonością
obliczeniową i~pełną interpretowalnością, co~czyni je~naturalnym punktem
odniesienia (ang.\ \emph{baseline}) przy~porównywaniu bardziej
zaawansowanych podejść~\cite{LibenNowell2007}.

Wstępna analiza struktury grafów DLG-DG-23 wykazała, że~miary oparte
na~wspólnych sąsiadach — Common Neighbors, indeks Jaccarda
i~miara Adamic-Adar — są nieskuteczne w~grafach dwudzielnych.
W~strukturze bipartite węzły tego samego typu (\emph{job}--\emph{job}
lub~\emph{table}--\emph{table}) nigdy nie~posiadają wspólnych sąsiadów
bezpośrednich, ponieważ krawędzie \emph{data\_flow} łączą wyłącznie
węzły różnych typów. W~konsekwencji wymienione miary zwracają wynik
zerowy dla~każdej typologicznie dopuszczalnej pary kandydackiej,
co~dyskwalifikuje je~jako metody wykrywania w~omawianym scenariuszu.
W~implementacji uwzględniono wyłącznie heurystyki działające poprawnie
na~grafach dwudzielnych: Preferential Attachment, L3, Katz
oraz~Personalized PageRank.

\subsubsection{Preferential Attachment}

Miara \emph{Preferential Attachment} (PA) opiera się na~obserwacji,
że~w~sieciach rzeczywistych nowe połączenia powstają z~prawdopodobieństwem
proporcjonalnym do~stopnia istniejących węzłów — węzły o~wielu sąsiadach
z~większym prawdopodobieństwem pozyskują kolejne
połączenia~\cite{Barabasi1999}. Jako heurystyka wykrywania krawędzi PA
została zastosowana przez Liben-Nowella i~Kleinberga~\cite{LibenNowell2007}.

Dla~pary węzłów $(u, v)$ wynik PA definiuje się jako:
%
\[
  \mathrm{PA}(u, v) = |N(u)| \cdot |N(v)|,
\]
%
gdzie $N(u)$ oraz $N(v)$ oznaczają zbiory sąsiadów odpowiednio węzła $u$
i~węzła $v$, a~$|\cdot|$ oznacza rozmiar zbioru, tj.\ stopień węzła w~grafie.

W~kontekście grafów lineage wysoki stopień węzła \emph{job} oznacza,
że~dane zadanie czyta z~wielu tabel źródłowych lub~zapisuje do~wielu
tabel docelowych; wysoki stopień węzła \emph{table} wskazuje na~tabelę
centralną — często pośrednią warstwę integracyjną w~potoku ETL.
Para $(u, v)$, w~której oba węzły mają duże stopnie, otrzymuje wysoki
wynik PA, co~interpretuje się jako sygnał potencjalnej zależności
\emph{data\_flow}. Miara PA nie~wymaga istnienia wspólnych sąsiadów,
co~stanowi jej kluczową przewagę w~grafach dwudzielnych.

\subsubsection{L3}

Miara L3 (ang.\ \emph{path of length 3}) liczy ścieżki długości~3
między węzłami $u$ i~$v$~\cite{LibenNowell2007}. W~grafie dwudzielnym
ścieżka długości~3 między węzłem \emph{job} a~węzłem \emph{table}
przebiega przez pośredni węzeł \emph{table} i~pośredni węzeł \emph{job}
— co odpowiada wzorcowi: zadanie~$u$ czyta z~tabeli~$t_1$, tabela~$t_1$
jest zapisywana przez zadanie~$j_1$, zadanie~$j_1$ czyta z~tabeli~$v$.
Wysoka wartość L3 sygnalizuje, że~para $(u, v)$ jest gęsto połączona
przez pośrednie obiekty, co~w~kontekście broken lineage może wskazywać
na~zależność ukrytą przez tabelę tymczasową.

Wynik L3 dla~pary $(u, v)$ definiuje się jako:
%
\[
  \mathrm{L3}(u, v) = \sum_{w \in V} A_{uw} \cdot A_{wv}^{(2)},
\]
%
gdzie $A$ oznacza macierz sąsiedztwa grafu, a~$A^{(2)} = A \cdot A$
— jej kwadrat (liczba ścieżek długości~2 między węzłami).
W~praktyce L3 obliczany jest jako element $(u,v)$ macierzy $A^3$.

\subsubsection{Katz}

Miara Katza uogólnia L3, sumując wkłady ścieżek wszystkich długości
z~wykładniczo malejącymi wagami~\cite{LibenNowell2007}:
%
\[
  \mathrm{Katz}(u, v) = \sum_{l=1}^{\infty} \beta^{l} \cdot |{\rm paths}_{uv}^{(l)}|
  = \left[(I - \beta A)^{-1} - I\right]_{uv},
\]
%
gdzie $\beta \in (0, 1)$ jest parametrem tłumienia, $A$ — macierzą
sąsiedztwa, a~$|{\rm paths}_{uv}^{(l)}|$ — liczbą ścieżek długości~$l$
między węzłami $u$ i~$v$. Mała wartość $\beta$ sprawia, że~dominujący
wkład pochodzi od~ścieżek najkrótszych. W~eksperymentach przyjęto
$\beta = 0{,}005$, co~zapewnia numeryczną stabilność obliczeń
przy~dużych grafach zbioru DLG-DG-23.

\subsubsection{Personalized PageRank}

\emph{Personalized PageRank} (PPR) realizuje ideę błądzenia losowego
z~restartem (\emph{random walk with restart})~\cite{LibenNowell2007}.
Dla~każdego węzła źródłowego $u$ oblicza się stacjonarny rozkład
prawdopodobieństwa błądzenia, które z~prawdopodobieństwem $\alpha$
powraca do~węzła $u$, a~z~prawdopodobieństwem $1 - \alpha$ przechodzi
do~losowego sąsiada. Wynik $\mathrm{PPR}(u, v)$ interpretuje się
jako prawdopodobieństwo odwiedzenia węzła $v$ podczas błądzenia
zainicjowanego w~$u$; wysoka wartość wskazuje na~bliskie sąsiedztwo
strukturalne. W~eksperymentach przyjęto $\alpha = 0{,}85$ (standardowa
wartość stosowana w~bibliotece NetworkX). Ze~względu na~koszt obliczenia
pełnej macierzy PPR dla~dużych grafów zastosowano obliczenia wsadowe
(\emph{batch computation}) z~wektorem startowym ograniczonym
do~węzłów zbioru treningowego.

% ============================================================
\subsection{Klasyczne metody uczenia maszynowego}
\label{sec:ml}
% ============================================================

Klasyczne metody uczenia maszynowego traktują problem wykrywania krawędzi
jako binarną klasyfikację, w~której wejściem jest wektor cech opisujący
parę węzłów $(u, v)$. Ze~względu na~anonimizację zbioru DLG-DG-23
(nazwy węzłów w~formacie UUID) jako cechy przyjęto wyłącznie atrybuty
topologiczne grafu — metryki oparte na~podobieństwie nazw (np.~Jaro-Winkler
stosowane w~\cite{Boinski2025}) są w~tym zbiorze nieaplikowalne.

\subsubsection{Wektor cech}

Dla~każdej pary kandydackiej $(u, v)$ wyznaczany jest wektor
11~cech topologicznych. Tabela~\ref{tab:cechy} przedstawia pełną
listę zastosowanych cech.

% TODO (WYMAGANE): Wstawić tabelę z opisem 11 cech:
%   Nr | Nazwa cechy | Opis | Uzasadnienie
%   Tabela musi być poprzedzona zdaniem „Tabela~X przedstawia...".
%   Uwzględnić: stopnie węzłów u i v, shortest_path, wyniki PA/L3/Katz/PPR,
%   przynależność do składowej spójnej, clustering coefficient itp.

Najwyższą istotność predykcyjną wykazuje cecha \emph{shortest\_path}
(długość najkrótszej ścieżki między $u$ a~$v$ w~grafie nieskierowanym),
co wskazuje, że~pary węzłów blisko połączonych w~istniejącej strukturze
grafu z~dużym prawdopodobieństwem są ze~sobą powiązane krawędzią
\emph{data\_flow} — nawet jeśli krawędź bezpośrednia nie~jest
obserwowana. Cecha ta~uzyskała ważność 52{,}6\% w~modelu Random Forest,
co~potwierdzono metodą permutacji na~zbiorze walidacyjnym.

\subsubsection{Random Forest}

Random Forest (RF) to~metoda uczenia zespołowego (\emph{ensemble
learning}) oparta na~uśrednianiu predykcji drzew decyzyjnych
wytrenowanych na~losowych podpróbkach danych i~cech
(\emph{bagging})~\cite{LibenNowell2007}. W~celu obsługi
niezrównoważonych klas zastosowano parametr
\texttt{class\_weight=\textquotedbl{}balanced\textquotedbl{}},
który skaluje wagi klas odwrotnie proporcjonalnie do~ich liczności.

\subsubsection{RUSBoost}

RUSBoost (\emph{Random Under-Sampling Boosting}) łączy algorytm
boostingu AdaBoost z~losowym podpróbkowaniem klasy większościowej
przed każdą iteracją. Metoda ta~jest szczególnie skuteczna przy silnym
niezbalansowaniu klas, typowym dla~problemu wykrywania krawędzi
w~rzadkich grafach. W~niniejszej pracy algorytm RUSBoost zaimplementowano
w~celu bezpośredniego porównania z~wynikami Boińskiego
et~al.~\cite{Boinski2025}, którzy uzyskali wartość F1~=~0{,}79
na~danych syntetycznych.

\subsubsection{LightGBM}

LightGBM to~implementacja gradientowego boostingu drzew decyzyjnych
(\emph{Gradient Boosted Decision Trees}) zoptymalizowana pod~kątem
wydajności i~skalowalności. W~odróżnieniu od~Random Forest, LightGBM
buduje kolejne drzewa w~sposób sekwencyjny, każde minimalizując błąd
resztkowy poprzednika. Parametr \texttt{scale\_pos\_weight} ustawiono
na~stosunek liczby próbek negatywnych do~pozytywnych, co~kompensuje
niezbalansowanie klas bez~konieczności podpróbowania.

% ============================================================
\subsection{Zanurzenia sieciowe — Node2Vec}
\label{sec:node2vec}
% ============================================================

\emph{Node2Vec} to~metoda uczenia reprezentacji węzłów grafu
(\emph{node embeddings}), opierająca się na~biased random walks
i~modelu Skip-gram~\cite{LibenNowell2007}. Każdemu węzłowi $v$
przypisywany jest wektor $\mathbf{z}_v \in \mathbb{R}^d$, zwany
zanurzeniem (\emph{embedding}), tak by~węzły o~podobnym kontekście
strukturalnym w~grafie były reprezentowane przez geometrycznie bliskie
wektory.

\subsubsection{Biased random walk}

Błądzenie losowe rozpoczyna się od~węzła $v$ i~na~każdym kroku
wybiera następny węzeł spośród sąsiadów bieżącego węzła
z~prawdopodobieństwem kontrolowanym przez dwa hiperparametry:

\begin{itemize}
  \item $p$ — parametr powrotu (\emph{return parameter}): kontroluje
        prawdopodobieństwo powrotu do~poprzednio odwiedzonego węzła;
        mała wartość $p$ sprzyja eksploatacji lokalnego sąsiedztwa,
  \item $q$ — parametr ekspansji (\emph{in-out parameter}): kontroluje
        skłonność do~oddalania się od~węzła startowego; wartość $q < 1$
        faworyzuje eksplorację dalszych obszarów grafu (strategia DFS),
        natomiast $q > 1$ — eksplorację lokalną (strategia BFS).
\end{itemize}

W~eksperymentach przyjęto $p = 1$ i~$q = 0{,}5$, co~sprzyja
eksploracji struktury globalnej grafu dwudzielnego — biased walk
z~$q < 1$ chętniej przemierza ścieżki prowadzące przez węzły obu
typów (\emph{job} i~\emph{table}), co~jest korzystne przy uczeniu
się reprezentacji par kandydackich.

\subsubsection{Generowanie embeddingu krawędzi}

Wykrywanie krawędzi wymaga przekształcenia pary zanurzeń węzłów
$(\mathbf{z}_u, \mathbf{z}_v)$ w~pojedynczy wektor reprezentujący
potencjalną krawędź $(u, v)$. Zastosowano iloczyn Hadamarda:
%
\[
  \mathbf{e}_{uv} = \mathbf{z}_u \odot \mathbf{z}_v,
\]
%
gdzie $\odot$ oznacza mnożenie element po~elemencie. Uzyskany
wektor $\mathbf{e}_{uv} \in \mathbb{R}^d$ podawany jest na~wejście
klasyfikatora — wielowarstwowej sieci neuronowej (ang.\
\emph{Multilayer Perceptron}, MLP) z~jedną warstwą ukrytą.

\subsubsection{Różnica względem grafowych sieci neuronowych}

Należy podkreślić, że~Node2Vec nie~jest grafową siecią neuronową
(ang.\ \emph{Graph Neural Network}, GNN). Metoda ta~uczy reprezentacji
węzłów wyłącznie na~podstawie sekwencji węzłów generowanych przez
błądzenie losowe (analogicznie do~metody Word2Vec na~sekwencjach
słów), bez~mechanizmu propagacji wiadomości (\emph{message passing})
charakterystycznego dla~GNN, takich jak GraphSAGE, R-GCN czy~HGT.
W~konsekwencji Node2Vec nie~uwzględnia atrybutów węzłów ani
typów krawędzi podczas uczenia zanurzeń. Rozszerzenie pracy
o~właściwe grafowe sieci neuronowe wskazano w~rozdziale~5 jako
kierunek dalszych badań.

% ============================================================

\chapter{Problem i~sposób jego rozwiązania}

Niniejszy rozdział precyzuje formalną definicję problemu broken lineage
i~jego ujęcie jako zadania klasyfikacji binarnej (sekcja~\ref{sec:sformulowanie}).
Sekcja~\ref{sec:architektura} opisuje architekturę zaproponowanego
potoku przetwarzania, a~sekcja~\ref{sec:dane} charakteryzuje zastosowany
zbiór danych. Metodykę symulacji brakujących krawędzi przedstawiono
w~sekcji~\ref{sec:scenariusze}, natomiast szczegóły implementacji
algorytmów omówiono w~sekcji~\ref{sec:algorytmy}.

% ============================================================
\section{Sformułowanie problemu}
\label{sec:sformulowanie}
% ============================================================

Graf lineage modeluje się jako graf skierowany $G^{*} = (V, E^{*})$,
gdzie $V$ oznacza zbiór obiektów bazodanowych (tabel, widoków, zadań ETL),
a~$E^{*}$ — zbiór wszystkich rzeczywistych relacji przepływu danych
(\emph{data\_flow}). W~praktyce część krawędzi jest niedostępna
z~powodu zjawiska broken lineage: tabele tymczasowe nie~są rejestrowane
w~standardowych słownikach bazy danych, a~funkcje definiowane przez
użytkownika (\emph{User-Defined Function}, UDF) ukrywają wewnętrzną
logikę transformacji. Obserwowany graf ma~zatem postać
$G_{\mathrm{obs}} = (V, E_{\mathrm{obs}})$, gdzie
$E_{\mathrm{obs}} \subsetneq E^{*}$. Brakujące krawędzie
$E^{*} \setminus E_{\mathrm{obs}}$ reprezentują zerwane zależności,
których odtworzenie jest celem niniejszej pracy.

Zadanie wykrywania brakujących krawędzi sprowadza się do~klasyfikacji
binarnej. Dla~każdej pary węzłów $(u, v) \in V \times V$ algorytm
wyznacza wartość $\hat{y}(u, v) \in [0, 1]$, interpretowaną jako
oszacowane prawdopodobieństwo istnienia krawędzi w~pełnym grafie
$G^{*}$. Para jest klasyfikowana jako krawędź istniejąca (klasa~1)
lub~nieistniejąca (klasa~0) po~przekroczeniu progu binaryzacji
wynoszącego $0{,}5$.

Ze~względu na~strukturę grafu DLG-DG-23 kandydatami do~wykrywania
są wyłącznie pary węzłów różnych typów: (\emph{job},~\emph{table})
lub~(\emph{table},~\emph{job}). Para (\emph{table},~\emph{table})
nigdy nie~jest połączona krawędzią \emph{data\_flow}, a~relacja
\emph{parent\_child} łączy wyłącznie węzły \emph{table} i~\emph{field}
— pary te nie~są przedmiotem wykrywania. Ograniczenie przestrzeni
kandydatów do~typologicznie dopuszczalnych par eliminuje trywialnie
fałszywe przykłady ujemne i~zapewnia spójność z~semantyką domeny.

% ============================================================
\section{Architektura rozwiązania}
\label{sec:architektura}
% ============================================================

Potok przetwarzania składa się z~czterech głównych komponentów:
modułu wczytywania danych (\emph{loader}), modułu podziału krawędzi
(\emph{splitter}), właściwych algorytmów wykrywania oraz modułu
ewaluacji.

% TODO (WYMAGANE): Wstawić rysunek blokowy potoku:
%   JSON → loader → NetworkX DiGraph → splitter → algorytmy → metryki.
%   Rysunek musi być wprowadzony zdaniem przed jego wystąpieniem,
%   np. „Rysunek~X przedstawia schemat blokowy zaproponowanego potoku."

\subsection{Moduł wczytywania danych}

Dane wejściowe dla~każdego grafu DLG przechowywane są w~dwóch plikach
JSON: \texttt{DLGi-node.json} zawiera opisy węzłów wraz z~atrybutem
\texttt{asset\_type}, natomiast \texttt{DLGi-edge.json} — opisy
krawędzi z~atrybutami \texttt{relation\_type} i~\texttt{relation\_id}.
Moduł \texttt{loader.py} wczytuje oba pliki dla~każdego grafu i~buduje
obiekt \texttt{nx.DiGraph} z~biblioteki NetworkX, a~następnie zwraca
słownik wszystkich 18~grafów zbioru DLG-DG-23. Sam graf przechowuje
atrybut \texttt{name} (np.\ \texttt{DLG1}).

\subsection{Moduł podziału krawędzi}

Moduł \texttt{splitter.py} wyodrębnia ze~struktury grafu krawędzie
typu \emph{data\_flow} i~dzieli je~na~trzy rozłączne podzbiory
w~proporcji 60/20/20 (treningowy, walidacyjny, testowy). Próbki
pozytywne stanowią krawędzie usunięte ze~zbioru treningowego
i~przeznaczone do~ewaluacji; próbki negatywne to~pary węzłów
typologicznie dopuszczalnych (wyłącznie \emph{job}--\emph{table}),
które nigdy nie~były połączone krawędzią. Stosunek próbek negatywnych
do~pozytywnych wynosi 1:1, co~zapewnia zbalansowany zbiór ewaluacyjny
i~stabilne wartości metryk porównawczych.

Moduł \texttt{scenario\_splitter.py} realizuje analogiczny podział,
lecz zamiast losowego usuwania krawędzi stosuje trzy scenariusze
symulacji broken lineage opisane w~sekcji~\ref{sec:scenariusze}.

\subsection{Moduł ewaluacji}

Moduł \texttt{metrics.py} wyznacza standardowe metryki klasyfikacji
binarnej: precyzję, czułość, miarę F1, AUC-ROC i~AUC-PR. Próg
binaryzacji wynosi $0{,}5$ dla~wszystkich metod. Wyniki są
uśredniane po~wszystkich 18~grafach zbioru DLG-DG-23 i~agregowane
według scenariusza oraz algorytmu, co~umożliwia bezpośrednie
porównanie między metodami.

% ============================================================
\section{Zbiór danych DLG-DG-23}
\label{sec:dane}
% ============================================================

Zastosowany zbiór danych pochodzi z~pracy Chen et~al.~\cite{Chen2024}
i~zawiera 18~grafów lineage pobranych ze~środowiska produkcyjnego
Huawei Cloud. Dane zostały zanonimizowane — nazwy tabel i~atrybutów
zastąpiono identyfikatorami w~formacie UUID, a~schematy bazy danych
nie~zostały udostępnione. Zbiór opublikowano w~2024~roku jako pierwszy
otwarty zestaw rzeczywistych grafów lineage; szczegółową charakterystykę
poszczególnych grafów zamieszczono w~tabeli~\ref{tab:dlg-stats} w~sekcji~\ref{sec:dane-opis}.

Zbiór DLG-DG-23 wybrano ze~względu na~jego unikatowy charakter:
jest to~jedyny publicznie dostępny zestaw rzeczywistych grafów lineage,
co~umożliwia ewaluację algorytmów na~danych o~autentycznej topologii
przemysłowej, bez~ryzyka artefaktów wprowadzonych przez dane
syntetyczne. Anonimizacja nazw węzłów ma~bezpośredni wpływ na~dobór
algorytmów: metryki podobieństwa oparte na~nazwach obiektów
(np.~odległość Jaro--Winkler stosowana przez Boińskiego
et~al.~\cite{Boinski2025}) są nieaplikowalne. W~niniejszej pracy
wektor cech ograniczono zatem do~atrybutów topologicznych grafu,
co~jednocześnie stanowi element badań nad~pytaniem, czy sama struktura
grafu lineage jest wystarczająca do~wykrywania zerwanych zależności.

\subsection{Charakterystyka grafów}
\label{sec:dane-opis}

Tabela~\ref{tab:dlg-stats} przedstawia statystyki poszczególnych
grafów zbioru DLG-DG-23.

\begin{table}[ht]
\centering
\caption{Charakterystyka grafów zbioru DLG-DG-23.}
\label{tab:dlg-stats}
\begin{threeparttable}
\begin{tabular}{llrl}
\toprule
\textbf{ID} & \textbf{Węzły} & \textbf{Krawędzie DATA\_FLOW} & \textbf{Skala} \\
\midrule
DLG1  &    298 &    24 & mały   \\
DLG2  &    464 &    37 & mały   \\
DLG3  &    603 &    57 & mały   \\
DLG4  &    415 &    32 & mały   \\
DLG5  & 1\,299 &   175 & średni \\
DLG6  & 13\,840 & 1\,443 & duży  \\
DLG7  &    574 &    43 & mały   \\
DLG8  &    546 &    34 & mały   \\
DLG9  &    756 &    62 & mały   \\
DLG10 & 5\,378 &   497 & duży   \\
DLG11 & 2\,024 &   121 & średni \\
DLG12 & 5\,645 &   469 & duży   \\
DLG13 & 2\,157 &   229 & średni \\
DLG14 &    453 &    24 & mały   \\
DLG15 &    558 &    41 & mały   \\
DLG16 & 5\,574 &   745 & duży   \\
DLG17 &    651 &    38 & mały   \\
DLG18 & 17\,085 & 1\,694 & duży \\
\bottomrule
\end{tabular}
\begin{tablenotes}
\small
\item Skala: mały $<$ 1\,000 węzłów; średni 1\,000–5\,000; duży $>$ 5\,000.
\item Liczba krawędzi DATA\_FLOW podana dla sumy zbiorów treningowego,
      walidacyjnego i~testowego (podział 60/20/20).
\end{tablenotes}
\end{threeparttable}
\end{table}

% ============================================================
\section{Metodyka symulacji broken lineage}
\label{sec:scenariusze}
% ============================================================

Ponieważ zbiór DLG-DG-23 zawiera w~pełni zaobserwowane grafy lineage,
zjawisko broken lineage symulowane jest poprzez kontrolowane usuwanie
krawędzi. Podejście to jest standardowe w~literaturze wykrywania
krawędzi~\cite{LibenNowell2007}: dysponując kompletnym grafem
$G^{*}$, usuwa się podzbiór krawędzi, tworząc graf obserwowany
$G_{\mathrm{obs}}$, a~następnie sprawdza, czy algorytm potrafi
odtworzyć usunięte krawędzie. W~niniejszej pracy przyjęto trzy
scenariusze symulacji, odzwierciedlające typowe przyczyny powstawania
broken lineage.

\textbf{Scenariusz A} — losowe usunięcie krawędzi — realizuje
najprostszą formę symulacji: 20\% krawędzi \emph{data\_flow} jest
usuwanych ze~zbioru treningowego w~sposób losowy, bez~względu na~typ
węzłów końcowych. Scenariusz ten modeluje sytuację, w~której~broken
lineage wynika z~przyczyn losowych (np.~niedokompletowania logów
systemowych) i~nie~jest skoncentrowany na~konkretnym typie obiektu.

\textbf{Scenariusz B} — usunięcie krawędzi tabeli pośredniej —
symuluje sytuację, w~której tabela pośrednia (staging) jest tabelą
tymczasową i~nie~zostaje zarejestrowana w~metadanych. Usuwane są
wszystkie krawędzie prowadzące do~wybranego węzła \emph{table}
o~wysokim stopniu (tabela pośrednia), co~powoduje odcięcie
całego podgrafu zależności przepływającego przez ten węzeł.

\textbf{Scenariusz C} — usunięcie krawędzi do~tabeli docelowej —
symuluje ukrycie zależności na~ostatnim etapie potoku ETL, kiedy
krawędź do~tabeli docelowej (\emph{data mart}) nie~zostaje utrwalona.
Usuwane są krawędzie wchodzące do~wybranego węzła \emph{table}
o~zerowym lub~niskim stopniu wyjściowym, co~odpowiada tabelom
końcowym w~hierarchii przepływu danych.

% TODO (WYMAGANE): Wstawić tabelę zbiorczą scenariuszy:
%   Scenariusz | Co usuwamy | Co symuluje | Typ węzłów dotkniętych
%   Tabela musi być poprzedzona zdaniem „Tabela X przedstawia...".

% TODO (OPCJONALNIE): Wstawić poglądowe rysunki małego grafu
%   przed i~po usunięciu krawędzi dla~każdego scenariusza.
%   Rysunki muszą być wprowadzone zdaniem w~tekście.

% ============================================================
\section{Zaimplementowane algorytmy}
\label{sec:algorytmy}
% ============================================================

W~pracy zaimplementowano trzy grupy algorytmów wykrywania
brakujących krawędzi. Pierwszą grupę stanowią heurystyki grafowe
— metody niewymagające uczenia, opierające się wyłącznie na~strukturze
grafu. Drugą grupę tworzą klasyczne metody uczenia maszynowego,
budujące klasyfikator na~podstawie wektora cech topologicznych
pary węzłów. Trzecia grupa obejmuje metody zanurzeń sieciowych
(\emph{network embeddings}), reprezentowane przez algorytm Node2Vec.

% ============================================================
\subsection{Heurystyki grafowe}
% ============================================================

Heurystyki grafowe stanowią najprostszą klasę metod wykrywania krawędzi
(\emph{link prediction}). Wynik dla~pary węzłów $(u, v)$ wyznaczany jest
wyłącznie na~podstawie lokalnej lub~globalnej struktury grafu, bez~etapu
trenowania modelu. Metody te charakteryzują się niską złożonością
obliczeniową i~pełną interpretowalnością, co~czyni je~naturalnym punktem
odniesienia (ang.\ \emph{baseline}) przy~porównywaniu bardziej
zaawansowanych podejść~\cite{LibenNowell2007}.

Wstępna analiza struktury grafów DLG-DG-23 wykazała, że~miary oparte
na~wspólnych sąsiadach — Common Neighbors, indeks Jaccarda
i~miara Adamic-Adar — są nieskuteczne w~grafach dwudzielnych.
W~strukturze bipartite węzły tego samego typu (\emph{job}--\emph{job}
lub~\emph{table}--\emph{table}) nigdy nie~posiadają wspólnych sąsiadów
bezpośrednich, ponieważ krawędzie \emph{data\_flow} łączą wyłącznie
węzły różnych typów. W~konsekwencji wymienione miary zwracają wynik
zerowy dla~każdej typologicznie dopuszczalnej pary kandydackiej,
co~dyskwalifikuje je~jako metody wykrywania w~omawianym scenariuszu.
W~implementacji uwzględniono wyłącznie heurystyki działające poprawnie
na~grafach dwudzielnych: Preferential Attachment, L3, Katz
oraz~Personalized PageRank.

\subsubsection{Preferential Attachment}

Miara \emph{Preferential Attachment} (PA) opiera się na~obserwacji,
że~w~sieciach rzeczywistych nowe połączenia powstają z~prawdopodobieństwem
proporcjonalnym do~stopnia istniejących węzłów — węzły o~wielu sąsiadach
z~większym prawdopodobieństwem pozyskują kolejne
połączenia~\cite{Barabasi1999}. Jako heurystyka wykrywania krawędzi PA
została zastosowana przez Liben-Nowella i~Kleinberga~\cite{LibenNowell2007}.

Dla~pary węzłów $(u, v)$ wynik PA definiuje się jako:
%
\[
  \mathrm{PA}(u, v) = |N(u)| \cdot |N(v)|,
\]
%
gdzie $N(u)$ oraz $N(v)$ oznaczają zbiory sąsiadów odpowiednio węzła $u$
i~węzła $v$, a~$|\cdot|$ oznacza rozmiar zbioru, tj.\ stopień węzła w~grafie.

W~kontekście grafów lineage wysoki stopień węzła \emph{job} oznacza,
że~dane zadanie czyta z~wielu tabel źródłowych lub~zapisuje do~wielu
tabel docelowych; wysoki stopień węzła \emph{table} wskazuje na~tabelę
centralną — często pośrednią warstwę integracyjną w~potoku ETL.
Para $(u, v)$, w~której oba węzły mają duże stopnie, otrzymuje wysoki
wynik PA, co~interpretuje się jako sygnał potencjalnej zależności
\emph{data\_flow}. Miara PA nie~wymaga istnienia wspólnych sąsiadów,
co~stanowi jej kluczową przewagę w~grafach dwudzielnych.

\subsubsection{L3}

Miara L3 (ang.\ \emph{path of length 3}) liczy ścieżki długości~3
między węzłami $u$ i~$v$~\cite{LibenNowell2007}. W~grafie dwudzielnym
ścieżka długości~3 między węzłem \emph{job} a~węzłem \emph{table}
przebiega przez pośredni węzeł \emph{table} i~pośredni węzeł \emph{job}
— co odpowiada wzorcowi: zadanie~$u$ czyta z~tabeli~$t_1$, tabela~$t_1$
jest zapisywana przez zadanie~$j_1$, zadanie~$j_1$ czyta z~tabeli~$v$.
Wysoka wartość L3 sygnalizuje, że~para $(u, v)$ jest gęsto połączona
przez pośrednie obiekty, co~w~kontekście broken lineage może wskazywać
na~zależność ukrytą przez tabelę tymczasową.

Wynik L3 dla~pary $(u, v)$ definiuje się jako:
%
\[
  \mathrm{L3}(u, v) = \sum_{w \in V} A_{uw} \cdot A_{wv}^{(2)},
\]
%
gdzie $A$ oznacza macierz sąsiedztwa grafu, a~$A^{(2)} = A \cdot A$
— jej kwadrat (liczba ścieżek długości~2 między węzłami).
W~praktyce L3 obliczany jest jako element $(u,v)$ macierzy $A^3$.

\subsubsection{Katz}

Miara Katza uogólnia L3, sumując wkłady ścieżek wszystkich długości
z~wykładniczo malejącymi wagami~\cite{LibenNowell2007}:
%
\[
  \mathrm{Katz}(u, v) = \sum_{l=1}^{\infty} \beta^{l} \cdot |{\rm paths}_{uv}^{(l)}|
  = \left[(I - \beta A)^{-1} - I\right]_{uv},
\]
%
gdzie $\beta \in (0, 1)$ jest parametrem tłumienia, $A$ — macierzą
sąsiedztwa, a~$|{\rm paths}_{uv}^{(l)}|$ — liczbą ścieżek długości~$l$
między węzłami $u$ i~$v$. Mała wartość $\beta$ sprawia, że~dominujący
wkład pochodzi od~ścieżek najkrótszych. W~eksperymentach przyjęto
$\beta = 0{,}005$, co~zapewnia numeryczną stabilność obliczeń
przy~dużych grafach zbioru DLG-DG-23.

\subsubsection{Personalized PageRank}

\emph{Personalized PageRank} (PPR) realizuje ideę błądzenia losowego
z~restartem (\emph{random walk with restart})~\cite{LibenNowell2007}.
Dla~każdego węzła źródłowego $u$ oblicza się stacjonarny rozkład
prawdopodobieństwa błądzenia, które z~prawdopodobieństwem $\alpha$
powraca do~węzła $u$, a~z~prawdopodobieństwem $1 - \alpha$ przechodzi
do~losowego sąsiada. Wynik $\mathrm{PPR}(u, v)$ interpretuje się
jako prawdopodobieństwo odwiedzenia węzła $v$ podczas błądzenia
zainicjowanego w~$u$; wysoka wartość wskazuje na~bliskie sąsiedztwo
strukturalne. W~eksperymentach przyjęto $\alpha = 0{,}85$ (standardowa
wartość stosowana w~bibliotece NetworkX). Ze~względu na~koszt obliczenia
pełnej macierzy PPR dla~dużych grafów zastosowano obliczenia wsadowe
(\emph{batch computation}) z~wektorem startowym ograniczonym
do~węzłów zbioru treningowego.

% ============================================================
\subsection{Klasyczne metody uczenia maszynowego}
\label{sec:ml}
% ============================================================

Klasyczne metody uczenia maszynowego traktują problem wykrywania krawędzi
jako binarną klasyfikację, w~której wejściem jest wektor cech opisujący
parę węzłów $(u, v)$. Ze~względu na~anonimizację zbioru DLG-DG-23
(nazwy węzłów w~formacie UUID) jako cechy przyjęto wyłącznie atrybuty
topologiczne grafu — metryki oparte na~podobieństwie nazw (np.~Jaro-Winkler
stosowane w~\cite{Boinski2025}) są w~tym zbiorze nieaplikowalne.

\subsubsection{Wektor cech}

Dla~każdej pary kandydackiej $(u, v)$ wyznaczany jest wektor
11~cech topologicznych. Tabela~\ref{tab:cechy} przedstawia pełną
listę zastosowanych cech.

% TODO (WYMAGANE): Wstawić tabelę z opisem 11 cech:
%   Nr | Nazwa cechy | Opis | Uzasadnienie
%   Tabela musi być poprzedzona zdaniem „Tabela~X przedstawia...".
%   Uwzględnić: stopnie węzłów u i v, shortest_path, wyniki PA/L3/Katz/PPR,
%   przynależność do składowej spójnej, clustering coefficient itp.

Najwyższą istotność predykcyjną wykazuje cecha \emph{shortest\_path}
(długość najkrótszej ścieżki między $u$ a~$v$ w~grafie nieskierowanym),
co wskazuje, że~pary węzłów blisko połączonych w~istniejącej strukturze
grafu z~dużym prawdopodobieństwem są ze~sobą powiązane krawędzią
\emph{data\_flow} — nawet jeśli krawędź bezpośrednia nie~jest
obserwowana. Cecha ta~uzyskała ważność 52{,}6\% w~modelu Random Forest,
co~potwierdzono metodą permutacji na~zbiorze walidacyjnym.

\subsubsection{Random Forest}

Random Forest (RF) to~metoda uczenia zespołowego (\emph{ensemble
learning}) oparta na~uśrednianiu predykcji drzew decyzyjnych
wytrenowanych na~losowych podpróbkach danych i~cech
(\emph{bagging})~\cite{LibenNowell2007}. W~celu obsługi
niezrównoważonych klas zastosowano parametr
\texttt{class\_weight=\textquotedbl{}balanced\textquotedbl{}},
który skaluje wagi klas odwrotnie proporcjonalnie do~ich liczności.

\subsubsection{RUSBoost}

RUSBoost (\emph{Random Under-Sampling Boosting}) łączy algorytm
boostingu AdaBoost z~losowym podpróbkowaniem klasy większościowej
przed każdą iteracją. Metoda ta~jest szczególnie skuteczna przy silnym
niezbalansowaniu klas, typowym dla~problemu wykrywania krawędzi
w~rzadkich grafach. W~niniejszej pracy algorytm RUSBoost zaimplementowano
w~celu bezpośredniego porównania z~wynikami Boińskiego
et~al.~\cite{Boinski2025}, którzy uzyskali wartość F1~=~0{,}79
na~danych syntetycznych.

\subsubsection{LightGBM}

LightGBM to~implementacja gradientowego boostingu drzew decyzyjnych
(\emph{Gradient Boosted Decision Trees}) zoptymalizowana pod~kątem
wydajności i~skalowalności. W~odróżnieniu od~Random Forest, LightGBM
buduje kolejne drzewa w~sposób sekwencyjny, każde minimalizując błąd
resztkowy poprzednika. Parametr \texttt{scale\_pos\_weight} ustawiono
na~stosunek liczby próbek negatywnych do~pozytywnych, co~kompensuje
niezbalansowanie klas bez~konieczności podpróbowania.

% ============================================================
\subsection{Zanurzenia sieciowe — Node2Vec}
\label{sec:node2vec}
% ============================================================

\emph{Node2Vec} to~metoda uczenia reprezentacji węzłów grafu
(\emph{node embeddings}), opierająca się na~biased random walks
i~modelu Skip-gram~\cite{LibenNowell2007}. Każdemu węzłowi $v$
przypisywany jest wektor $\mathbf{z}_v \in \mathbb{R}^d$, zwany
zanurzeniem (\emph{embedding}), tak by~węzły o~podobnym kontekście
strukturalnym w~grafie były reprezentowane przez geometrycznie bliskie
wektory.

\subsubsection{Biased random walk}

Błądzenie losowe rozpoczyna się od~węzła $v$ i~na~każdym kroku
wybiera następny węzeł spośród sąsiadów bieżącego węzła
z~prawdopodobieństwem kontrolowanym przez dwa hiperparametry:

\begin{itemize}
  \item $p$ — parametr powrotu (\emph{return parameter}): kontroluje
        prawdopodobieństwo powrotu do~poprzednio odwiedzonego węzła;
        mała wartość $p$ sprzyja eksploatacji lokalnego sąsiedztwa,
  \item $q$ — parametr ekspansji (\emph{in-out parameter}): kontroluje
        skłonność do~oddalania się od~węzła startowego; wartość $q < 1$
        faworyzuje eksplorację dalszych obszarów grafu (strategia DFS),
        natomiast $q > 1$ — eksplorację lokalną (strategia BFS).
\end{itemize}

W~eksperymentach przyjęto $p = 1$ i~$q = 0{,}5$, co~sprzyja
eksploracji struktury globalnej grafu dwudzielnego — biased walk
z~$q < 1$ chętniej przemierza ścieżki prowadzące przez węzły obu
typów (\emph{job} i~\emph{table}), co~jest korzystne przy uczeniu
się reprezentacji par kandydackich.

\subsubsection{Generowanie embeddingu krawędzi}

Wykrywanie krawędzi wymaga przekształcenia pary zanurzeń węzłów
$(\mathbf{z}_u, \mathbf{z}_v)$ w~pojedynczy wektor reprezentujący
potencjalną krawędź $(u, v)$. Zastosowano iloczyn Hadamarda:
%
\[
  \mathbf{e}_{uv} = \mathbf{z}_u \odot \mathbf{z}_v,
\]
%
gdzie $\odot$ oznacza mnożenie element po~elemencie. Uzyskany
wektor $\mathbf{e}_{uv} \in \mathbb{R}^d$ podawany jest na~wejście
klasyfikatora — wielowarstwowej sieci neuronowej (ang.\
\emph{Multilayer Perceptron}, MLP) z~jedną warstwą ukrytą.

\subsubsection{Różnica względem grafowych sieci neuronowych}

Należy podkreślić, że~Node2Vec nie~jest grafową siecią neuronową
(ang.\ \emph{Graph Neural Network}, GNN). Metoda ta~uczy reprezentacji
węzłów wyłącznie na~podstawie sekwencji węzłów generowanych przez
błądzenie losowe (analogicznie do~metody Word2Vec na~sekwencjach
słów), bez~mechanizmu propagacji wiadomości (\emph{message passing})
charakterystycznego dla~GNN, takich jak GraphSAGE, R-GCN czy~HGT.
W~konsekwencji Node2Vec nie~uwzględnia atrybutów węzłów ani
typów krawędzi podczas uczenia zanurzeń. Rozszerzenie pracy
o~właściwe grafowe sieci neuronowe wskazano w~rozdziale~5 jako
kierunek dalszych badań.
